% Fichier principal du rapport.
%
% Moteur : pdflatex, par latexmk. Le choix est mesuré et consigné à
% docs/decisions/ : babel en français et les polices Latin Modern sont présents
% dans toute distribution complète, là où xelatex et lualatex exigeraient des
% polices du système, qui diffèrent entre la machine de développement et
% l'exécuteur d'intégration continue. La compilation qui fait foi est celle de
% l'intégration continue.
%
% Aucun contenu n'est rédigé dans ce fichier : il n'assemble.

\documentclass[12pt,a4paper]{report}

\usepackage[T1]{fontenc}
\usepackage[utf8]{inputenc}
\usepackage{lmodern}
\usepackage[french]{babel}
\usepackage{csquotes}

\usepackage[a4paper,margin=2.5cm]{geometry}
\usepackage{microtype}

% L'aération du texte : interligne à 1,15 et espacement de paragraphe visible, sans retrait.
% Le document était composé en interligne simple et sans blanc entre paragraphes, ce qui le
% rendait dense à lire. Le contenu ne change pas ; la page en porte moins.
%
% `setspace` appartient à l ensemble `collection-latexrecommended`, le même que `microtype` et `booktabs`
% chargés ci-dessus. L'image de composition de l'intégration continue les compose déjà à chaque
% exécution verte : le paquet y est donc nécessairement présent, et le choix est mesuré plutôt
% que supposé.
\usepackage{setspace}
\setstretch{1.15}
\setlength{\parskip}{0.6em}
\setlength{\parindent}{0pt}
\usepackage{longtable}
\usepackage{booktabs}
\usepackage{fancyhdr}

% L'inclusion d'images. Le paquet appartient à l'ensemble `graphics`, requis par
% LaTeX 2e et présent dans toute distribution, et le choix est MESURÉ et non supposé de deux
% façons : `kpsewhich graphicx.sty` le trouve ici sous
% `texmf-dist/tex/latex/graphics/graphicx.sty`, et surtout `pgfcore.sty` — que
% `pgfplots`, déjà chargé ci-dessous, tire à lui — porte à sa ligne 10
% `\RequirePackage{graphicx}`. L'image de composition de l'intégration continue
% compose donc déjà ce paquet à chaque exécution verte, sans que rien ne l'ait
% demandé explicitement.
\usepackage{graphicx}

% La couleur, pour les intitulés de la page de garde. Elle est déjà chargée par la chaîne de
% `pgfplots` — le journal de compilation nomme `xcolor.sty` — mais un paquet dont on se sert se
% charge explicitement plutôt que par effet de bord. Même ensemble `collection-latexrecommended` que
% `microtype` et `setspace`.
\usepackage{xcolor}
\definecolor{griscartouche}{gray}{0.42}

% Tracé vectoriel. Le paquet est présent dans toute distribution complète, et l'image de
% composition de l'intégration continue en porte une : le choix est mesuré, non supposé.
\usepackage{pgfplots}
\pgfplotsset{compat=1.18}

% Les nombres des graduations et des tableaux composés suivent l'usage français : virgule
% décimale, espace fine insécable aux milliers. Sans cette ligne, pgfplots écrit « 1,000 » là où
% le rendu du registre des chiffres écrit « 1 000 », et deux conventions typographiques
% coexisteraient dans la même page.
\pgfplotsset{/pgf/number format/.cd, use comma, 1000 sep={\,}}

% La mise en tableau d'un fichier de données, du même paquet — mesuré : `tlmgr search --file
% --global pgfplotstable.sty` rend `pgfplots: texmf-dist/tex/latex/pgfplots/pgfplotstable.sty`.
% Aucune distribution qui porte l'un ne peut manquer l'autre, et c'est ce qui permet à un tableau
% de chiffres de venir d'une commande plutôt que d'être tapé.
\usepackage{pgfplotstable}

% ------------------------------------------------------------------------------------------
% La correspondance entre relations injectées et conclusions.
%
% Trois commandes, et un contrôle les lit DANS LES DEUX SENS : toute relation du registre des
% relations injectées est reprise par une conclusion ou déclarée non reprise avec son motif, et
% toute conclusion renvoie à une relation ou déclare qu'elle n'en vient pas. Une ligne retirée
% d'un côté ou de l'autre fait rougir ce contrôle en nommant l'orpheline.
% ------------------------------------------------------------------------------------------
% ------------------------------------------------------------------------------------------
% L'annonce du plan, dans l'introduction générale.
%
% Un titre de chapitre annoncé s'écrit `\annonceChapitre{numéro}{titre}` et se compose en
% italique, au fil de la prose. Un contrôle lit ces appels et les confronte AUX FICHIERS DE
% CHAPITRE, dans les deux sens : tout chapitre annoncé existe sous ce numéro avec ce titre, et
% tout chapitre existant est annoncé. Sans lui, une renumérotation laisserait l'introduction
% annoncer un plan qui n'est plus celui du document — le projet a déjà eu, à l'écran, un
% décompte écrit à la main devenu faux, et il ne l'a pas trouvé par un contrôle.
% ------------------------------------------------------------------------------------------
\newcommand{\annonceChapitre}[2]{\emph{#2}}

\newcommand{\conclusion}[1]{\textbf{#1.}}
\newcommand{\relreprise}[4]{#1 & #2 & #3 & #4 \\}
\newcommand{\relnonreprise}[2]{#1 & --- & \multicolumn{2}{p{0.55\linewidth}}{#2} \\}
\newcommand{\conclusionhors}[2]{#1 & #2 \\}

\usepackage[backend=biber,style=numeric,sorting=none]{biblatex}
\addbibresource{biblio.bib}

% Les adresses en ligne des entrées bibliographiques débordaient dans la marge — mesuré, jusqu'à
% 333 pt, soit près de douze centimètres. Ces trois pénalités sont celles que `biblatex` expose
% pour autoriser une coupure d'adresse ; sans elles, une adresse longue ne se coupe jamais.
\setcounter{biburlnumpenalty}{7000}
\setcounter{biburlucpenalty}{7000}
\setcounter{biburllcpenalty}{7000}

\usepackage[hidelinks]{hyperref}

% Les adresses en ligne de la bibliographie débordaient dans la marge, jusqu'à 333 pt — près de
% douze centimètres. Les pénalités de `biblatex` ne suffisent pas : elles n'autorisent la coupure
% qu'entre certaines classes de caractères, et ces adresses n'en offrent pas. `xurl` autorise la
% coupure partout dans une adresse. Il est du même ensemble `collection-latexextra` que `csquotes`,
% chargé plus haut et composé à chaque exécution verte de l'intégration continue.
\usepackage{xurl}

% La hauteur d'en-tête : `fancyhdr` avertissait à chaque compilation qu'elle était trop petite, et
% l'avertissement était présent depuis l'origine. La valeur ci-dessous est celle que le paquet
% demande une fois l'interligne élargi — mesurée sur son propre message, non devinée.
\setlength{\headheight}{15.5pt}
\addtolength{\topmargin}{-3.5pt}

% Une marge d'élasticité de dernier recours. Sans elle, une ligne que TeX ne sait pas couper
% déborde dans la marge ; avec elle, il desserre l'interlettrage de cette ligne-là plutôt que de
% la laisser dépasser. Elle ne peut rien contre un mot plus large que la colonne — une empreinte
% de soixante-quatre caractères, par exemple.
\setlength{\emergencystretch}{2em}

\pagestyle{fancy}
\fancyhf{}
\fancyhead[L]{\leftmark}
\fancyfoot[C]{\thepage}

% Les marqueurs nominatifs du rapport, et le seul endroit du dépôt où un nom de
% personne peut figurer. Tout autre fichier les emploie par leur commande, jamais
% en écrivant la valeur.
%
% Un marqueur laissé vide fait échouer un contrôle : c'est la propriété que le
% critère de terminaison du projet exige, et tests/test_marqueurs_nominatifs.py
% la vérifie sur CE fichier, jamais sur le PDF produit.
%
% Renseigner une valeur, c'est écrire entre les accolades de la commande. Ne pas
% renommer les commandes : le contrôle les cherche par leur nom.

% --- l'état du document -------------------------------------------------------
% Deux valeurs, et deux seulement.
%
%   brouillon : le rapport n'est pas remis. Les cinq marqueurs nominatifs
%               ci-dessous sont TOUS vides, et la page de garde porte une mention
%               visible de version de travail.
%   remise    : le rapport est remis. Les cinq marqueurs sont TOUS renseignés, et
%               la mention disparaît.
%
% Un contrôle vérifie l'accord entre cet état et les cinq marqueurs : renseigner
% quatre marqueurs sur cinq est rouge dans les deux états, et c'est l'oubli qui
% arrive. Passer à la remise ne coûte que ce mot — et le contrôle dira alors quel
% marqueur manque encore.
\newcommand{\etatDuDocument}{brouillon}

% --- marqueurs nominatifs : ils portent un nom de personne ---------------------
\newcommand{\marqueurAuteur}{}
\newcommand{\marqueurEncadrantAcademique}{}
\newcommand{\marqueurEncadrantProfessionnel}{}
\newcommand{\marqueurPresidentJury}{}
\newcommand{\marqueurExaminateur}{}

% --- marqueurs de contexte : ils ne portent aucun nom de personne --------------
% Le contrôle ne les exige pas renseignés ; ils sont ici pour que la page de
% garde n'écrive aucune valeur en dur.
\newcommand{\marqueurEtablissementFormation}{}
\newcommand{\marqueurFiliere}{}
\newcommand{\marqueurAnneeUniversitaire}{}
\newcommand{\marqueurOrganismeAccueil}{Hôpital Sidi Saïd, Meknès}
\newcommand{\marqueurDateSoutenance}{}

% --- ce que l'état fait porter au document ------------------------------------
% Le mécanisme est typographique : le document se suffit, et ne dépend d'aucun
% contrôle pour dire ce qu'il est. `\siBrouillon{...}` compose son argument en
% état de brouillon, et rien en état de remise.
%
% La comparaison passe par \ifx sur deux commandes développées : c'est la forme
% que le noyau offre sans paquet supplémentaire, et elle compare les définitions
% caractère pour caractère.
\newcommand{\etatBrouillonAttendu}{brouillon}
\newcommand{\siBrouillon}[1]{%
  \ifx\etatDuDocument\etatBrouillonAttendu#1\fi%
}

% La mention elle-même, écrite une fois et employée par les deux documents.
\newcommand{\mentionVersionDeTravail}{%
  \siBrouillon{\textbf{Version de travail} — document non remis}%
}

% --- les trois étiquettes de provenance -------------------------------------------------------
% La prose porte les mêmes étiquettes que les colonnes de l'entrepôt, et pour la même raison :
% un lecteur doit pouvoir dire, sur n'importe quelle affirmation, d'où elle vient.
%
%   DOC  une entrée du registre des sources l'étaye. Elle se cite par \cite{...}, et la clé
%        renvoie au fichier bibliographique produit à partir du registre.
%   OBS  un relevé d'observation l'étaye. Aucune source publiée ne la porte : elle a été vue.
%   HYP  aucune des deux. C'est une convention posée, et le rapport la déclare comme telle.
%
% Chaque fichier de chapitre déclare en tête l'ensemble de ce sur quoi il repose, sous ces trois
% étiquettes, et un contrôle vérifie que la déclaration coïncide exactement avec ce que le fichier
% porte — dans les deux sens.

% --- ces deux étiquettes NE COMPOSENT PLUS RIEN ------------------------------------------------
% Elles portaient un exposant `[obs.]` et `[conv.]` au fil de la prose. Elles ne l'impriment plus,
% et le motif est le DESTINATAIRE du document : ce rapport s'adresse à un encadrant de stage, non
% à une relecture adverse. Un exposant tous les deux paragraphes est un appareil de défense, et il
% se lit comme tel.
%
% L'APPAREIL RESTE ENTIÈREMENT CÂBLÉ, IL CESSE SEULEMENT D'IMPRIMER. L'argument reste l'identifiant
% du relevé ou de la convention ; chaque fichier de chapitre continue de déclarer en tête ce sur
% quoi il repose ; et `tests/test_provenance_des_chapitres.py` continue de confronter la
% déclaration à ce que le fichier porte, dans les deux sens. CE CONTRÔLE LIT LES SOURCES DE
% COMPOSITION, JAMAIS LE PDF : rien de ce qu'il vérifie ne dépend de ce que ces deux commandes
% composent, et les rendre muettes ne l'affaiblit d'aucune façon. Retirer un `\releve{...}` d'une
% phrase reste rouge, exactement comme avant.
%
% Les citations `\cite{...}` de l'étiquette DOC, elles, restent visibles : renvoyer à une source
% publiée est l'usage académique normal, et ce n'est pas un appareil de défense.

% Un relevé d'observation. L'argument est son identifiant, court et stable ; il n'est pas composé.
\newcommand{\releve}[1]{}

% Une convention posée. Même forme, même règle.
\newcommand{\convention}[1]{}

% --- ce qui reste à rédiger --------------------------------------------------------------------
% Un paragraphe qui relève de la rédaction personnelle et que le squelette ne peut pas écrire.
% Le premier argument est son identifiant d'observation, le second ce qu'il portera.
%
% La marque est VISIBLE à la compilation, toujours : un rapport remis en la portant se voit à
% l'œil. Le contrôle, lui, s'appuie sur l'état déclaré du document — il rougit en état de remise
% et s'abstient en état de brouillon, où un paragraphe non encore écrit est normal.
\newcommand{\aRediger}[2]{%
  \par\medskip\noindent\fbox{%
    \begin{minipage}{0.95\linewidth}%
      \textbf{À rédiger} — #2%
    \end{minipage}%
  }\par\medskip%
}

% Les emplacements d'image du rapport, et les seules commandes par lesquelles une
% image entre dans le document.
%
% DEUX FAMILLES D'IMAGE SONT ADMISES, ET DEUX SEULEMENT : le logotype de
% l'établissement de formation, et les captures du tableau de bord PRODUIT PAR CE
% PROJET. Toute image du système d'information observé reste interdite, sans
% exception — c'est la décision `0010`, que rien ici n'assouplit.
% `tests/test_aucune_image.py` tient la propriété sur les fichiers suivis : seuls
% `report/figures/logos/` et `report/figures/tableau-de-bord/` la portent.
%
% LA PLACE EST RÉSERVÉE MAINTENANT ET REMPLIE PLUS TARD. Chaque commande teste
% l'existence du fichier attendu. S'il existe, elle l'inclut. S'il n'existe pas,
% elle compose un cadre de la MÊME hauteur ou de la même largeur déclarée, portant
% la légende et la mention « à insérer ». Le document compose donc à tous les
% stades, et la pagination ne bouge presque pas à l'insertion — c'est tout
% l'intérêt du mécanisme, et une commande qui composerait un cadre d'une autre
% taille que l'image finale l'annulerait.
%
% L'ÉCART RÉSIDUEL EST MESURÉ, PAS SUPPOSÉ, ET LA MESURE A CORRIGÉ LE CODE. Le
% cadre est posé avec `\fboxsep` nul et aligné PAR LE BAS — `\parbox[b]`, comme
% `\includegraphics`, dont la ligne de base est le bas de l'image. Un premier
% essai alignait le cadre sur son centre (`\parbox[c]`) : un document témoin
% composé avec puis sans l'image montrait alors la ligne suivante déplacée de
% 3,304 pt. Avec l'alignement par le bas, la même mesure donne 0,395 pt — la seule
% épaisseur du trait inférieur du cadre, soit 0,047 % de la hauteur utile d'une
% page A4. Un emplacement rempli déplace donc le texte qui le suit d'un vingtième
% de ligne, et la pagination ne bouge pas.
%
% La mesure : composer un document portant `\captureTdb` et un mot repérable
% après lui, relever `yMin` de ce mot par `pdftotext -bbox`, avec le fichier image
% en place puis retiré. Les deux relevés sont au rapport de ce travail.

% Le trait du cadre d'attente, et l'absence de marge intérieure. Les deux sont
% posés localement par chaque commande, jamais globalement : `\fbox` sert ailleurs
% dans le document — la mention de version de travail en porte un — et lui retirer
% sa marge partout changerait cette mention-là aussi.
\newcommand{\cadreDAttenteReglages}{\setlength{\fboxsep}{0pt}\setlength{\fboxrule}{0.4pt}}

% Deux longueurs de travail. Les dimensions ne se calculent pas par juxtaposition
% dans un argument de macro : `4#1` avec `#1` valant `2.2cm` donnerait `42.2cm`, et
% `0.625#2` produirait une dimension inexistante. Le calcul passe donc par
% `\setlength`, qui admet un facteur devant une longueur.
\newlength{\largeurEmplacement}
\newlength{\hauteurEmplacement}

% Un cadre d'attente de dimensions imposées. La largeur et la hauteur sont celles
% que l'image occupera ; le texte est centré dans les deux sens et n'influe pas
% sur la boîte, `\parbox[c][hauteur][c]` fixant la hauteur indépendamment de son
% contenu.
\newcommand{\cadreDAttente}[3]{%
  {\cadreDAttenteReglages%
   \fbox{\parbox[b][#2][c]{#1}{\centering\footnotesize\itshape #3\par}}}%
}

% ------------------------------------------------------------------------------
% LES DEUX LOGOTYPES INSTITUTIONNELS.
%
%   \logosInstitutionnels
%
% La page de garde en porte DEUX, côte à côte : l'établissement de formation et
% l'organisme dont il relève. La commande précédente, `\logoEcole{hauteur}`, n'en
% prenait qu'un et cherchait un fichier `ecole.pdf` ou `ecole.png` qui n'a jamais
% existé ; elle est remplacée, et non doublée, pour qu'il n'y ait qu'une seule
% voie par laquelle un logotype entre dans le document.
%
% LES DEUX HAUTEURS NE SONT PAS ÉGALES, ET C'EST MESURÉ. Imposer la même hauteur
% aux deux ne les équilibre pas : un logotype large et bas paraît plus grand qu'un
% logotype compact et haut à hauteur égale. L'égalisation porte donc sur la
% SURFACE OPTIQUE — la surface de la boîte utile, marge transparente exclue.
%
%   fichier          dimensions   rapport   boîte utile   marge   encre        hauteur
%   INSEA-logo.png   1200x1304    0,9202    983x1110      30,3 %  582 421 px   20,0 mm
%   HCP-logo.png     3790x1088    3,4835    3781x1088      0,2 %  547 086 px   17,2 mm
%
% Le calcul : la surface portée à la page vaut la surface en pixels multipliée par
% le carré de l'échelle, et l'échelle vaut la hauteur déclarée divisée par la
% hauteur du fichier. Égaliser donne un rapport d'échelles égal à la racine du
% rapport des surfaces, et le rapport des hauteurs s'en déduit en multipliant par
% le rapport des hauteurs en pixels.
%
% DEUX SURFACES SONT CANDIDATES, ET ELLES NE DONNENT PAS LE MÊME RÉSULTAT. La
% boîte utile donne un rapport de hauteurs de 2,3267, donc 8,6 mm pour le second
% logotype ; l'encre — le nombre de pixels réellement opaques — donne 1,1616, donc
% 17,2 mm. L'écart tient à ceci : la boîte utile de l'INSEA est couverte à 53,4 %,
% celle du HCP à 13,3 % seulement, parce que ce second logotype est un mot écrit
% en fin, largement étalé, quand le premier est une marque pleine. La boîte utile
% créditait donc le HCP d'une surface qu'il n'occupe pas.
%
% LA MESURE À L'ENCRE EST RETENUE, ET LE RENDU L'A TRANCHÉE. La page composée avec
% 8,6 mm a été regardée en image : le second logotype y paraissait nettement plus
% faible que le premier, et sa ligne « Haut-Commissariat au Plan » n'était plus
% lisible. À 17,2 mm les deux marques pèsent le même poids à l'œil et les deux
% textes se lisent. C'est la surface d'encre qui mesure la présence optique, pas
% le rectangle qui l'entoure.
%
% LES DEUX SONT ALIGNÉS SUR LEUR LIGNE MÉDIANE, jamais sur leur base : deux
% logotypes de hauteurs différentes alignés par le bas paraissent décalés.
% `\raisebox{-0.5\height}` pose chaque boîte à cheval sur la ligne de base, si
% bien que les deux milieux se retrouvent à la même hauteur. L'encombrement
% vertical de la ligne reste celui du plus grand des deux, et il ne dépend pas de
% la présence des fichiers : les cadres d'attente ont les mêmes dimensions.
%
% CENTRER LES BOÎTES NE CENTRE PAS L'ENCRE, et la mesure l'a montré. Le fichier de
% l'INSEA porte 51 pixels de marge transparente en haut et 143 en bas : son encre
% est donc plus haute que le centre de sa boîte de 46 pixels, soit 0,706 mm à la
% hauteur déclarée. Le fichier du HCP, lui, n'a aucune marge verticale. Les deux
% boîtes alignées, les deux encres l'étaient à 0,76 mm près — relevé sur l'image du
% PDF, et cohérent avec le calcul. Chaque logotype porte donc une CORRECTION,
% mesurée sur son fichier, qui ramène sa ligne médiane d'encre sur la ligne de
% base. La correction s'applique au `\raisebox` extérieur, donc aussi bien à
% l'image qu'à son cadre d'attente.
\newcommand{\hauteurLogoEtablissement}{20mm}
\newcommand{\hauteurLogoTutelle}{17.2mm}
\newcommand{\ecartEntreLogos}{16mm}

% Un logotype, ou son cadre d'attente aux mêmes dimensions.
%
%   \unLogo{base du nom}{hauteur}{rapport largeur/hauteur}{correction}{légende}
%
% Le rapport est celui du FICHIER, mesuré sur lui, et non une valeur commode : il
% ne sert qu'au cadre d'attente, dont la largeur doit être celle qu'occupera
% l'image pour que retirer un fichier ne déplace rien, ni verticalement ni
% horizontalement. La correction est la longueur dont il faut descendre la boîte
% pour que sa médiane d'encre tombe sur la ligne de base ; elle vaut zéro pour un
% fichier sans marge transparente.
\newcommand{\unLogo}[5]{%
  \raisebox{\dimexpr-0.5\height-#4\relax}{%
    \IfFileExists{figures/logos/#1.pdf}{%
      \includegraphics[height=#2]{figures/logos/#1.pdf}%
    }{%
      \IfFileExists{figures/logos/#1.png}{%
        \includegraphics[height=#2]{figures/logos/#1.png}%
      }{%
        \setlength{\hauteurEmplacement}{#2}%
        \setlength{\largeurEmplacement}{#3\hauteurEmplacement}%
        \cadreDAttente{\largeurEmplacement}{\hauteurEmplacement}{#5\\à insérer}%
      }%
    }%
  }%
}

\newcommand{\logosInstitutionnels}{%
  \unLogo{INSEA-logo}{\hauteurLogoEtablissement}{0.9202}{0.706mm}{logotype de l'établissement}%
  \hspace{\ecartEntreLogos}%
  \unLogo{HCP-logo}{\hauteurLogoTutelle}{3.4835}{0mm}{logotype de la tutelle}%
}

% ------------------------------------------------------------------------------
% Une capture du tableau de bord produit par ce projet.
%
%   \captureTdb{fichier}{largeur}{légende}{label}
%
% `fichier` est le nom seul, sans répertoire ni extension : la commande le cherche
% sous `figures/tableau-de-bord/`, en `.png` puis en `.pdf`. Écrire le répertoire à
% chaque appel l'exposerait à diverger du préfixe que le contrôle tolère.
%
% LA HAUTEUR RÉSERVÉE VIENT DE LA LARGEUR, PAR UN RAPPORT DÉCLARÉ PAR CAPTURE.
% Une capture n'a pas de hauteur déclarée à l'appel, et le cadre d'attente doit
% pourtant en avoir une. Un premier état retenait un rapport UNIQUE de 16:10 pour
% toutes les captures, en écrivant que c'était le point aveugle du mécanisme :
% « une capture d'un autre rapport se compose correctement, mais déplace ce qui la
% suit ». Les trois captures déposées ont démenti la supposition — 2,456, 1,973 et
% 1,964 au lieu de 1,600 —, et le rapport unique réservait donc de 1,7 à 3,2 cm de
% trop pour chacune. Le rapport est désormais DÉCLARÉ POUR CHAQUE FICHIER, mesuré
% sur le fichier lui-même ; la valeur unique ne subsiste que comme défaut, pour une
% capture non encore déposée et donc non encore mesurable.
%
% La déclaration est ici, dans le fichier qui tient les images, et non aux appels :
% c'est une propriété du fichier image, pas du texte qui l'appelle.
\newcommand{\rapportHauteurCapture}{0.625}

\newcommand{\declarerCapture}[2]{%
  \expandafter\gdef\csname rapportCapture@#1\endcsname{#2}%
}
\newcommand{\rapportDeLaCapture}[1]{%
  \ifcsname rapportCapture@#1\endcsname
    \csname rapportCapture@#1\endcsname
  \else
    \rapportHauteurCapture
  \fi
}

% Hauteur divisée par largeur, en pixels, de chaque fichier déposé.
\declarerCapture{page-activite}{0.40724}%  1374/3374
\declarerCapture{page-donnees}{0.50680}%   1714/3382
\declarerCapture{page-sejours}{0.50918}%   1720/3378

% La numérotation. Le rapport ne porte aujourd'hui aucun environnement flottant et
% aucune `\caption` ; un `\label` posé sans compteur renverrait à la section
% courante, ce qui ferait mentir tout `\ref`. Un compteur propre est donc déclaré,
% et `\refstepcounter` en fait la cible du `\label` qui suit.
\newcounter{capturetdb}
\renewcommand{\thecapturetdb}{\arabic{capturetdb}}

\newcommand{\legendeCapture}[2]{%
  \begin{center}%
    \footnotesize Figure~\thecapturetdb~— #1\label{#2}%
  \end{center}%
}

% L'IMAGE EST POSÉE DANS LA BOÎTE RÉSERVÉE, ET NON À CÔTÉ. La boîte a la hauteur
% déclarée pour ce fichier ; l'image y entre bornée en largeur ET en hauteur, à
% rapport conservé. Un fichier dont le rapport dériverait de sa déclaration ne
% déborderait donc pas de la réservation : il y laisserait un peu de blanc. C'est
% ce qui rend la pagination indépendante de la présence du fichier, au lieu de la
% faire dépendre d'une supposition écrite ailleurs.
\newcommand{\captureTdb}[4]{%
  \refstepcounter{capturetdb}%
  \setlength{\largeurEmplacement}{#2}%
  \edef\facteurDeLaCapture{\rapportDeLaCapture{#1}}%
  \setlength{\hauteurEmplacement}{\facteurDeLaCapture\largeurEmplacement}%
  \begin{center}%
    \IfFileExists{figures/tableau-de-bord/#1.png}{%
      \parbox[b][\hauteurEmplacement][c]{\largeurEmplacement}{\centering%
        \includegraphics[width=\largeurEmplacement,height=\hauteurEmplacement,keepaspectratio]{figures/tableau-de-bord/#1.png}}%
    }{%
      \IfFileExists{figures/tableau-de-bord/#1.pdf}{%
        \parbox[b][\hauteurEmplacement][c]{\largeurEmplacement}{\centering%
          \includegraphics[width=\largeurEmplacement,height=\hauteurEmplacement,keepaspectratio]{figures/tableau-de-bord/#1.pdf}}%
      }{%
        \cadreDAttente{\largeurEmplacement}{\hauteurEmplacement}{#3\\capture à insérer}%
      }%
    }%
  \end{center}%
  \legendeCapture{#3}{#4}%
}

% Fichier produit mécaniquement à partir de docs/chiffres/registre_chiffres.yml.
% Ne pas modifier à la main : docs/chiffres/generer_chiffres_tex.py le réécrit.
%
% Le rapport n'écrit aucune valeur en clair. Il écrit \chiffre{identifiant}, et la valeur vient
% d'ici, mesurée par la commande que le registre porte en regard de cet identifiant.
%
% Un identifiant inconnu déclenche une erreur de compilation qui le nomme : un chiffre absent doit
% arrêter la composition, jamais laisser un blanc dans une phrase.

\newcommand{\chiffre}[1]{%
  \ifcsname chiffre@#1\endcsname
    \csname chiffre@#1\endcsname
  \else
    \GenericError{}{Chiffre inconnu : #1}{}{Cet identifiant n'existe pas au registre des chiffres.}%
  \fi
}


\expandafter\def\csname chiffre@source-patients-lignes\endcsname{29\,107}
\expandafter\def\csname chiffre@source-rendez-vous-lignes\endcsname{14\,169}
\expandafter\def\csname chiffre@source-passages-lignes\endcsname{40\,650}
\expandafter\def\csname chiffre@source-passages-urgences-lignes\endcsname{27\,360}
\expandafter\def\csname chiffre@source-mouvements-lignes\endcsname{6\,310}
\expandafter\def\csname chiffre@source-prises-en-charge-lignes\endcsname{16\,430}
\expandafter\def\csname chiffre@source-factures-lignes\endcsname{21\,066}
\expandafter\def\csname chiffre@source-lignes-facture-lignes\endcsname{160\,936}
\expandafter\def\csname chiffre@source-encaissements-lignes\endcsname{23\,231}
\expandafter\def\csname chiffre@source-creances-lignes\endcsname{5\,876}
\expandafter\def\csname chiffre@source-relances-lignes\endcsname{1\,014}
\expandafter\def\csname chiffre@source-lignes-total\endcsname{346\,149}
\expandafter\def\csname chiffre@source-tables\endcsname{11}
\expandafter\def\csname chiffre@quarantaine-patients-lignes\endcsname{64}
\expandafter\def\csname chiffre@quarantaine-rendez-vous-lignes\endcsname{42}
\expandafter\def\csname chiffre@quarantaine-lignes-total\endcsname{106}
\expandafter\def\csname chiffre@colonnes-source-registre\endcsname{175}
\expandafter\def\csname chiffre@colonnes-source-obs\endcsname{81}
\expandafter\def\csname chiffre@colonnes-source-doc\endcsname{72}
\expandafter\def\csname chiffre@colonnes-source-hyp\endcsname{22}
\expandafter\def\csname chiffre@colonnes-source-catalogue\endcsname{175}
\expandafter\def\csname chiffre@colonnes-source-catalogue-obs\endcsname{81}
\expandafter\def\csname chiffre@colonnes-source-catalogue-doc\endcsname{72}
\expandafter\def\csname chiffre@colonnes-source-catalogue-hyp\endcsname{22}
\expandafter\def\csname chiffre@colonnes-aval-total\endcsname{398}
\expandafter\def\csname chiffre@colonnes-aval-decrites\endcsname{398}
\expandafter\def\csname chiffre@colonnes-intermediate\endcsname{175}
\expandafter\def\csname chiffre@colonnes-marts\endcsname{223}
\expandafter\def\csname chiffre@parametres-total\endcsname{256}
\expandafter\def\csname chiffre@parametres-fichiers\endcsname{25}
\expandafter\def\csname chiffre@parametres-obs\endcsname{5}
\expandafter\def\csname chiffre@parametres-doc\endcsname{52}
\expandafter\def\csname chiffre@parametres-hyp\endcsname{199}
\expandafter\def\csname chiffre@capacite-fonctionnelle\endcsname{40}
\expandafter\def\csname chiffre@journees-2024\endcsname{7\,851}
\expandafter\def\csname chiffre@admissions-2024\endcsname{1\,197}
\expandafter\def\csname chiffre@consultations-2024\endcsname{4\,142}
\expandafter\def\csname chiffre@prelevements-2024\endcsname{9\,625}
\expandafter\def\csname chiffre@examens-2024\endcsname{49\,597}
\expandafter\def\csname chiffre@interventions-2024\endcsname{0}
\expandafter\def\csname chiffre@tom-publie\endcsname{53,8}
\expandafter\def\csname chiffre@dms-publie\endcsname{6,6}
\expandafter\def\csname chiffre@periode-jours\endcsname{912}
\expandafter\def\csname chiffre@periode-debut\endcsname{2024-01-01}
\expandafter\def\csname chiffre@periode-fin\endcsname{2026-06-30}
\expandafter\def\csname chiffre@sejours-produits\endcsname{2\,980}
\expandafter\def\csname chiffre@consultations-produites\endcsname{10\,310}
\expandafter\def\csname chiffre@passages-urgences-produits\endcsname{27\,360}
\expandafter\def\csname chiffre@patients-produits\endcsname{29\,107}
\expandafter\def\csname chiffre@personnes-distinctes\endcsname{25\,842}
\expandafter\def\csname chiffre@urgences-scenario-retenu\endcsname{30}
\expandafter\def\csname chiffre@urgences-borne-basse\endcsname{14}
\expandafter\def\csname chiffre@urgences-scenario-median\endcsname{30}
\expandafter\def\csname chiffre@urgences-borne-haute\endcsname{54}
\expandafter\def\csname chiffre@defaut-taux-doublons\endcsname{0,04}
\expandafter\def\csname chiffre@defaut-dates-aberrantes\endcsname{0,003}
\expandafter\def\csname chiffre@defaut-ages-incoherents\endcsname{0,005}
\expandafter\def\csname chiffre@defaut-rdv-doublon-creneau\endcsname{0,012}
\expandafter\def\csname chiffre@defaut-factures-sans-pec\endcsname{0,03}
\expandafter\def\csname chiffre@defaut-familles\endcsname{13}
\expandafter\def\csname chiffre@part-sexe-masculin\endcsname{0,497}
\expandafter\def\csname chiffre@part-urbain\endcsname{0,639}
\expandafter\def\csname chiffre@part-sans-couverture\endcsname{0,2}
\expandafter\def\csname chiffre@part-couverture-principale\endcsname{0,5}
\expandafter\def\csname chiffre@tranches-age\endcsname{6}
\expandafter\def\csname chiffre@niveaux-tri\endcsname{5}
\expandafter\def\csname chiffre@part-tri-niveau-4\endcsname{0,41}
\expandafter\def\csname chiffre@part-tri-niveau-5\endcsname{0,22}
\expandafter\def\csname chiffre@part-sejours-depuis-urgences\endcsname{0,57}
\expandafter\def\csname chiffre@mouvements-lits-distincts\endcsname{40}
\expandafter\def\csname chiffre@factures-avec-creance\endcsname{5\,486}


\begin{document}

% Page de garde. Aucune image du système d'information, ici comme ailleurs ; le
% seul logotype admis est celui de l'établissement de formation, et il vient de
% `\logoEcole`, qui compose un cadre de la même hauteur tant que le fichier
% n'existe pas.
%
% Aucune valeur n'est écrite ici : chaque champ vient d'un marqueur.
%
% AUCUN TEXTE FIXE N'ENTOURE UN MARQUEUR QUI PEUT ÊTRE VIDE. Chaque champ dont le
% marqueur peut ne pas être renseigné passe par `\siRenseigne` : l'étiquette
% disparaît avec sa valeur au lieu d'annoncer une phrase qui s'arrête. La règle
% vaut pour les deux noms, pour la date de soutenance, et pour les trois champs de
% contexte. Le seul champ inconditionnel est l'organisme d'accueil, dont le
% marqueur est renseigné dans le fichier des marqueurs lui-même — et il tient la
% table des noms non vide, ce qui évite un tableau sans aucune ligne.

\begin{titlepage}
\centering
\vspace*{1cm}

\logoEcole{2.2cm}

\vspace{0.8cm}

\siRenseigne{\marqueurEtablissementFormation}{%
  {\large \marqueurEtablissementFormation\par}\vspace{0.5cm}%
}
\siRenseigne{\marqueurFiliere}{{\large \marqueurFiliere\par}}

\vfill

\siBrouillon{{\large\fbox{\mentionVersionDeTravail}\par}\vspace{1cm}}

{\Huge\bfseries Rapport de stage d'application\par}
\vspace{1cm}
{\Large Chaîne de données et tableau de bord\\pour un établissement hospitalier\par}

\vfill

\begin{tabular}{rl}
  \siRenseigne{\marqueurAuteur}{Réalisé par & \marqueurAuteur \\[0.4em]}
  \siRenseigne{\marqueurEncadrantProfessionnel}{Encadrant de stage & \marqueurEncadrantProfessionnel \\[0.4em]}
  Organisme d'accueil & \marqueurOrganismeAccueil \\
\end{tabular}

\vfill

\siRenseigne{\marqueurDateSoutenance}{%
  {\large Soutenu le \marqueurDateSoutenance\par}\vspace{0.5cm}%
}
\siRenseigne{\marqueurAnneeUniversitaire}{%
  {\large Année universitaire \marqueurAnneeUniversitaire\par}%
}

\end{titlepage}


\pagenumbering{roman}
% Les remerciements.
%
% AUCUN NOM DE PERSONNE EN CLAIR. L'encadrant de stage est nommé par son marqueur, comme sur la page
% de garde, et le marqueur reste vide tant que `noms.tex` n'est pas déposé. La phrase se compose
% alors sans nom, ce qui est le comportement attendu en état de brouillon.
%
% Rien n'est affirmé qui ne soit établi. Le seul fait dont ce dossier porte la trace est que
% l'encadrant a remis le cahier de charges et laissé toute liberté d'action dans le cadre qu'il
% fixe : c'est ce que les remerciements disent, et rien de plus.

\chapter*{Remerciements}
\addcontentsline{toc}{chapter}{Remerciements}

Je remercie mon encadrant de
stage\siRenseigne{\marqueurEncadrantProfessionnel}{, \marqueurEncadrantProfessionnel,} pour la
confiance qu'il m'a accordée. Il m'a remis le cahier de charges, en a précisé le contenu, et m'a
laissé toute liberté sur la démarche à suivre dans le cadre qu'il avait fixé. Cette liberté a
déplacé vers moi la responsabilité du cadrage, et c'est ce qui a rendu ce travail formateur.

Je remercie le service d'accueil et d'admission qui a hébergé ce travail, à
l'\marqueurOrganismeAccueil{}, pour son accueil et pour le temps qu'il m'a consacré au poste
d'observation.

Je remercie enfin\siRenseigne{\marqueurEtablissementFormation}{ l'\marqueurEtablissementFormation{},}
pour la formation qui m'a permis d'aborder ce travail.

\chapter*{Résumé}
\addcontentsline{toc}{chapter}{Résumé}

\section*{Mots-clés}

\chapter*{Abstract}
\addcontentsline{toc}{chapter}{Abstract}

\section*{Keywords}


\tableofcontents

\cleardoublepage
\pagenumbering{arabic}

% Ce chapitre repose sur :
%   DOC: (aucun)
%   OBS: (aucun)
%   HYP: (aucun)

\chapter*{Introduction générale}
\addcontentsline{toc}{chapter}{Introduction générale}


\chapter{L'organisme d'accueil}

\section{Le système de santé marocain et la région Fès-Meknès}

\section{Le centre hospitalier préfectoral de Meknès}

\section{L'hôpital Sidi Saïd}

\section{Le service d'accueil et le périmètre d'observation}

% Ce chapitre repose sur :
%   DOC: s06, s07, s08, s09, s10
%   OBS: REL-DEM.I01, REL-DEM.I02, REL-DEM.P01, REL-DEM.P02, REL-DEM.P03, REL-DEM.P04, REL-DEM.P05,
%   OBS: REL-DEM.M01, REL-DEM.M02, REL-DEM.M03, REL-DEM.M04, REL-IPP.O01, REL-IPP.O02, REL-IPP.O03,
%   OBS: REL-IPP.O04, REL-IPP.C01, REL-IPP.C02, REL-IPP.C03, REL-IPP.C04, REL-IPP.C05, REL-IPP.C06,
%   OBS: REL-IPP.C07, REL-IPP.C08, REL-IPP.C09, REL-IPP.C10, REL-IPP.C11, REL-IPP.C12, REL-IPP.C13,
%   OBS: REL-IPP.F01, REL-IPP.F02, REL-IPP.F03, REL-IPP.F04, REL-IPP.R01, REL-IPP.R02, REL-IPP.R03,
%   OBS: REL-IPP.R04, REL-IPP.R05, REL-IPP.R06, REL-IPP.R07, REL-IPP.R08, REL-IPP.R09, REL-IPP.R10,
%   OBS: REL-IPP.R11, REL-IPP.R12, REL-PAT.T01, REL-PAT.T02, REL-PAT.T03, REL-PAT.D01, REL-PAT.D02,
%   OBS: REL-PAT.D03, REL-PAT.D04, REL-PAT.D05, REL-PAT.D06, REL-PAT.D07, REL-PAT.D08, REL-PAT.D09,
%   OBS: REL-PAT.D10, REL-PAT.D11, REL-PAT.D12, REL-PAT.D13, REL-PAT.D14, REL-PAT.D15, REL-PAT.D16,
%   OBS: REL-PAT.A01, REL-PAT.A02, REL-PAT.A03, REL-PAT.A04, REL-PAT.A05, REL-PAT.A06, REL-PAT.H01,
%   OBS: REL-PAT.H02, REL-PAT.H03, REL-PAT.H04, REL-PAT.H05, REL-PAT.H06, REL-PAT.H07, REL-PAT.H08,
%   OBS: REL-PAT.H09, REL-PAT.H10, REL-PAT.H11, REL-PAT.H12, REL-PAT.H13, REL-PAT.H14, REL-PAT.H15,
%   OBS: REL-PAT.N01, REL-PAT.N02, REL-PAT.N03, REL-PAT.N04, REL-PAT.N05, REL-PAT.N06, REL-PAT.K01,
%   OBS: REL-RDV.I01, REL-RDV.I02, REL-RDV.I03, REL-RDV.I04, REL-RDV.I05, REL-RDV.I06, REL-RDV.I07,
%   OBS: REL-RDV.I08, REL-RDV.I09, REL-RDV.I10, REL-RDV.R01, REL-RDV.R02, REL-RDV.R03, REL-RDV.R04,
%   OBS: REL-RDV.R05, REL-RDV.R06, REL-RDV.R07, REL-RDV.R08, REL-RDV.R09, REL-RDV.R10, REL-RDV.R11,
%   OBS: REL-RDV.R12, REL-RDV.R13, REL-RDV.R14, REL-RDV.C01, REL-RDV.C02, REL-RDV.C03, REL-RDV.C04,
%   OBS: REL-RDV.C05, REL-RDV.C06, REL-RDV.C07, REL-RDV.C08, REL-RDV.L01, REL-RDV.L02, REL-RDV.L03
%   HYP: (aucun)
%   CHI: (aucun)
%   SER: (aucune)
%   SECTIONS: 8

\chapter{Le système d'information hospitalier}

\section{Hosix.NET, origine et architecture}

Le système d'information en place est Hosix.NET, édité par SIVSA Soluciones Informáticas
\cite{s06}. Sa conception remonte aux années 1980, et l'éditeur en revendique le déploiement dans
plus de trente hôpitaux répartis sur six pays \cite{s07}.

L'origine du produit se lit dans sa nomenclature avant de se lire dans sa documentation
commerciale. La raison sociale de l'éditeur est espagnole, et les modules portent des noms
espagnols que la documentation reprend sans les traduire : \emph{Consultas} pour la gestion des
agendas et des rendez-vous \cite{s08}, \emph{Facturación} pour l'émission des factures
\cite{s09}, auxquels s'ajoutent \emph{Médicos}, \emph{Urgencias} et \emph{admisión}. Un produit
déployé au Maroc et servi en français conserve donc, dans son organisation interne, la langue de
son concepteur.

L'architecture fonctionnelle se contracte en trois paquets — Core, Clinic et Extension Packs —
qui rassemblent vingt-six modules : douze au Core, six au Clinic, huit aux Extension Packs
\cite{s06}. La répartition n'est pas indifférente. Le Core porte les fonctions administratives du
dossier — admission, agendas, facturation —, le Clinic les fonctions médicales, et les Extension
Packs des compléments contractables séparément. Un établissement peut donc exploiter la gestion
administrative des patients sans exploiter la couche clinique, et c'est une configuration que
l'observation conduite ici ne permet ni de confirmer ni d'exclure : le périmètre installé n'est
pas connu, la documentation de l'éditeur décrivant ce qui est contractable et jamais ce qui est
contracté.

Le produit a fait l'objet d'une migration technique, de WebForms vers le patron
modèle-vue-contrôleur, et l'éditeur la déclare achevée sur quatre modules : prescription,
admission, facturation et approvisionnements \cite{s07}. Cette migration n'est pas un détail
d'implantation. Elle explique qu'un même logiciel présente des écrans dont l'ergonomie diffère
d'un module à l'autre, et elle situe l'observation conduite ici dans un produit en cours de
transformation, non dans un état stable.

Une précision sur ce que ce chapitre ne dit pas, et elle vaut pour tout le dossier
documentaire. La documentation consultée est celle de l'éditeur, et elle est commerciale : elle
publie une fiche par module, et ces fiches énoncent des fonctions, pas des schémas de données.
Elles établissent qu'une fonction existe et ce qu'elle recouvre ; elles ne disent jamais quels
champs la portent, sous quel type, ni sous quelle contrainte. Aucune colonne du modèle construit
plus loin ne peut donc être justifiée par la seule documentation de l'éditeur.

Tout ce qui touche aux champs vient du relevé conduit au poste. C'est une matière d'une autre
nature : elle est datée, limitée à un profil et à quatre écrans, et elle ne couvre qu'une fraction
du produit. Elle a en revanche une propriété que la documentation n'a pas — elle porte sur
l'installation réelle de l'établissement, et non sur le produit générique. Les quatre sections de
relevé qui suivent la restituent intégralement, élément par élément.

\section{Les cinq profils applicatifs et les processus qu'ils portent}

L'écran de démarrage propose cinq profils applicatifs, et le relevé les reproduit sous leur
libellé exact : \texttt{MSM - FACTURATION SAA}, \texttt{MSM - GESTION DE RDV},
\texttt{MSM - RECOUVREMENT}, \texttt{MSM - URGENCE} et
\texttt{MSM RAPPORTS ET STATISTIQUES}\releve{REL-DEM.P02}. L'irrégularité du dernier — l'absence
de tiret séparateur — est celle de l'écran, et elle est conservée : un libellé se reproduit, il ne
se corrige pas.

Ces cinq profils ne sont pas cinq vues d'un même écran : chacun ouvre un jeu de modules distinct,
et les habilitations cloisonnent l'accès. Un seul a été ouvert au poste, les quatre autres n'ont
été relevés qu'à l'écran de démarrage. Le tableau ci-dessous rattache chaque profil au processus
qu'il porte et au module de l'éditeur qui le sert, en distinguant ce qui a été vu de ce qui a été
reconstruit par voie documentaire.

\begin{center}
\begin{tabular}{@{}p{4.3cm}p{4.6cm}p{3.1cm}p{2.2cm}@{}}
\toprule
Profil & Processus porté & Module & Établi par \\
\midrule
\texttt{MSM - GESTION DE RDV} & prise, modification et annulation de rendez-vous ; liste d'attente & \emph{Consultas} \cite{s08} & observation \\
\addlinespace
\texttt{MSM - FACTURATION SAA} & captation des mouvements, émission des factures & \emph{Facturación} \cite{s09} & documentation \\
\addlinespace
\texttt{MSM - RECOUVREMENT} & suivi des créances et des relances & \emph{Facturación} \cite{s09} & documentation \\
\addlinespace
\texttt{MSM - URGENCE} & admission et prise en charge aux urgences & \emph{Urgencias} & documentation \\
\addlinespace
\texttt{MSM RAPPORTS ET STATISTIQUES} & production des indicateurs remontés à la statistique nationale & — \cite{s10} & documentation \\
\bottomrule
\end{tabular}
\end{center}

\medskip

La documentation du module \emph{Consultas} énonce ce que le profil de rendez-vous recouvre :
génération d'agendas, prise de rendez-vous manuelle ou automatique, modification, annulation,
intégration à la liste d'attente, production de notifications, de listings et de statistiques
\cite{s08}. Celle de \emph{Facturación} énonce la captation manuelle ou automatique des
mouvements, le contrôle de l'état facturé ou en attente, le paramétrage par groupes, par types de
contrats et de prestations, et la facturation aux entités publiques, privées et mutuelles
\cite{s09}. L'éditeur ne publie aucune fiche distincte pour le recouvrement : la fonction est
intégrée au module de facturation, ce qui explique que deux profils s'y rattachent.

Le cinquième profil est le seul dont le contenu se déduise d'une source extérieure à l'éditeur :
la nomenclature des indicateurs qu'un hôpital nommé remonte à la statistique nationale
\cite{s10}.

\section{Méthode de relevé}

L'observation a porté sur un poste de travail du service, sur le profil
\texttt{MSM - GESTION DE RDV}, le 28 juillet 2026\releve{REL-DEM.P02}. Quatre écrans ont été
parcourus : l'écran de démarrage, la recherche d'identifiant patient, la fiche patient et l'écran
de prise de rendez-vous. Cent quarante-trois éléments y ont été relevés.

Chaque élément est consigné avec un identifiant stable, le bloc de l'écran auquel il appartient,
son libellé reproduit à l'identique, son type apparent, et la valeur affichée lorsqu'une valeur
l'était. L'identifiant n'est pas une commodité de rédaction : c'est la clé par laquelle le
registre de provenance des colonnes désigne son témoin d'observation, et par laquelle les
tableaux qui suivent se relient au modèle de données construit plus loin.

Deux conventions gouvernent la restitution, et toutes deux consistent à ne pas affirmer plus que
ce qui a été vu. La première : le caractère obligatoire ou facultatif d'un champ n'a été relevé
sur aucun écran, et n'est donc affirmé nulle part — une valeur binaire imposée ici produirait cent
quarante-trois affirmations non observées. La seconde : la mention \texttt{sans objet} désigne un
élément qui n'est pas un champ de saisie de dossier — bouton, onglet, entrée de menu, libellé de
profil, colonne de résultat, filtre de recherche. Les boutons des barres d'outils ne figurent pas
aux tableaux : ils déclenchent une opération et ne portent aucune donnée, et le relevé les déclare
comme tels.

\section{Relevé — l'écran de démarrage et les profils}

L'écran d'entrée identifie l'établissement et l'autorité dont il relève, puis propose les cinq
profils et quatre menus principaux. Onze éléments y sont relevés.

\begin{longtable}{@{}p{2.9cm}p{3.9cm}p{2.8cm}p{4.0cm}@{}}
\toprule
Identifiant & Libellé à l'écran & Type apparent & Valeur observée \\
\midrule\endfirsthead
\toprule
Identifiant & Libellé à l'écran & Type apparent & Valeur observée \\
\midrule\endhead
\multicolumn{4}{@{}l}{\textbf{Identification du site}} \\
\texttt{REL-DEM.I01} & Identification du site & texte & \texttt{005 Fès-Meknès; Sidi Said} \\
\texttt{REL-DEM.I02} & Ministère de tutelle & texte & \texttt{Ministère de la Santé} \\
\addlinespace
\multicolumn{4}{@{}l}{\textbf{Profils applicatifs}} \\
\texttt{REL-DEM.P01} & MSM - FACTURATION SAA & liste & — \\
\texttt{REL-DEM.P02} & MSM - GESTION DE RDV & liste & — \\
\texttt{REL-DEM.P03} & MSM - RECOUVREMENT & liste & — \\
\texttt{REL-DEM.P04} & MSM - URGENCE & liste & — \\
\texttt{REL-DEM.P05} & MSM RAPPORTS ET STATISTIQUES & liste & — \\
\addlinespace
\multicolumn{4}{@{}l}{\textbf{Menus principaux}} \\
\texttt{REL-DEM.M01} & Hosix.NET & menu & — \\
\texttt{REL-DEM.M02} & Données Patient & menu & — \\
\texttt{REL-DEM.M03} & Centre de Consultation & menu & — \\
\texttt{REL-DEM.M04} & Factures & menu & — \\
\addlinespace
\bottomrule
\end{longtable}

\medskip

Le code de région y est accolé au nom de l'établissement sous une graphie que le relevé
conserve telle quelle, séparateur compris\releve{REL-DEM.I01}. Les quatre menus donnent l'ossature
de la navigation, et deux d'entre eux mènent aux écrans relevés ci-après.

\section{Relevé — la recherche d'identifiant patient}

Cet écran est celui par lequel un agent retrouve un dossier avant toute opération. Trente-trois
éléments y sont relevés, répartis en quatre blocs.

\begin{longtable}{@{}p{2.9cm}p{3.9cm}p{2.8cm}p{4.0cm}@{}}
\toprule
Identifiant & Libellé à l'écran & Type apparent & Valeur observée \\
\midrule\endfirsthead
\toprule
Identifiant & Libellé à l'écran & Type apparent & Valeur observée \\
\midrule\endhead
\multicolumn{4}{@{}l}{\textbf{Onglets}} \\
\texttt{REL-IPP.O01} & Chercher & onglet & — \\
\texttt{REL-IPP.O02} & Chercher MPI & onglet & — \\
\texttt{REL-IPP.O03} & Chercher MPI Pondéré & onglet & — \\
\texttt{REL-IPP.O04} & Nouvel IPP & onglet & — \\
\addlinespace
\multicolumn{4}{@{}l}{\textbf{Critères de recherche}} \\
\texttt{REL-IPP.C01} & N° document & texte & \texttt{C.I.N} \\
\texttt{REL-IPP.C02} & Nom & texte & — \\
\texttt{REL-IPP.C03} & Prénom & texte & — \\
\texttt{REL-IPP.C04} & Date de naissance & date & — \\
\texttt{REL-IPP.C05} & Mois & entier & — \\
\texttt{REL-IPP.C06} & Année & entier & — \\
\texttt{REL-IPP.C07} & Approximation en années & entier & — \\
\texttt{REL-IPP.C08} & Âge & entier & — \\
\texttt{REL-IPP.C09} & Téléphone & texte & — \\
\texttt{REL-IPP.C10} & N° document de contact & texte & — \\
\texttt{REL-IPP.C11} & Sexe & liste & — \\
\texttt{REL-IPP.C12} & État civil & liste & — \\
\texttt{REL-IPP.C13} & N° assuré & texte & — \\
\addlinespace
\multicolumn{4}{@{}l}{\textbf{Filtres de population}} \\
\texttt{REL-IPP.F01} & Patients Hospitalisés & option & — \\
\texttt{REL-IPP.F02} & Patients aux Urgences & option & — \\
\texttt{REL-IPP.F03} & Tous les patients & option & — \\
\texttt{REL-IPP.F04} & Exitus & case\ a\ cocher & — \\
\addlinespace
\multicolumn{4}{@{}l}{\textbf{Colonnes de résultat}} \\
\texttt{REL-IPP.R01} & IPP & colonne & — \\
\texttt{REL-IPP.R02} & Nom & colonne & — \\
\texttt{REL-IPP.R03} & D. Naissance & colonne & — \\
\texttt{REL-IPP.R04} & N° pièce d'identité & colonne & — \\
\texttt{REL-IPP.R05} & N° Assu & colonne & — \\
\texttt{REL-IPP.R06} & E. Civil & colonne & — \\
\texttt{REL-IPP.R07} & Tél & colonne & — \\
\texttt{REL-IPP.R08} & Sexe & colonne & — \\
\texttt{REL-IPP.R09} & Province & colonne & — \\
\texttt{REL-IPP.R10} & Ville & colonne & — \\
\texttt{REL-IPP.R11} & Code postal & colonne & — \\
\texttt{REL-IPP.R12} & Probabilité & colonne & — \\
\addlinespace
\bottomrule
\end{longtable}

\medskip

\textbf{Un index maître des patients existe dans le système.} Le bloc d'onglets en porte la
preuve : à côté de la recherche ordinaire et de la création d'un nouvel identifiant, deux onglets
s'adressent explicitement à un index maître, dont l'un en recherche
pondérée\releve{REL-IPP.O03}. Le bloc des colonnes de résultat le confirme : la dernière colonne
affiche une \emph{probabilité}\releve{REL-IPP.R12}. Une recherche qui rend des probabilités n'est
pas une recherche exacte — c'est un rapprochement.

Ce que le relevé établit s'arrête là. Il atteste qu'un index maître et une recherche pondérée
existent dans l'installation observée ; il n'établit rien sur leur déploiement au-delà de cet
établissement, et aucune source du dossier ne le porte. La portée de cette observation pour la
modélisation est reprise au chapitre~\ref{chap:qualite}.

Les treize critères de recherche méritent une lecture. Ils ne cherchent pas seulement par
identifiant : nom, prénom, date de naissance, mais aussi mois et année seuls, âge, et une
\emph{approximation en années}. Un écran qui offre de chercher à l'âge approché est un écran conçu
pour des dossiers dont la date de naissance est incertaine ou absente.

\section{Relevé — la fiche patient}

La fiche patient est l'écran le plus dense du périmètre observé : cinquante-sept éléments, dont
dix boutons de barre d'outils qui ne figurent pas au tableau. Les quarante-sept restants se
répartissent en trois onglets et cinq blocs de données.

\begin{center}
\begin{tabular}{@{}c@{}}
\fbox{\begin{minipage}{0.86\linewidth}
\centering\small
\fbox{\begin{minipage}{0.95\linewidth}\centering Barre d'actions \\ \footnotesize dix boutons, hors tableau\end{minipage}}\\[3pt]
\fbox{\begin{minipage}{0.95\linewidth}\centering Onglets : Général \quad|\quad Contacts \quad|\quad Quota soins primaires\end{minipage}}\\[3pt]
\fbox{\begin{minipage}{0.95\linewidth}\centering Données principales \\ \footnotesize seize éléments\end{minipage}}\\[3pt]
\fbox{\begin{minipage}{0.46\linewidth}\centering Compagnie d'assurance \\ \footnotesize six éléments\end{minipage}}%
\hfill
\fbox{\begin{minipage}{0.46\linewidth}\centering Né \\ \footnotesize six éléments\end{minipage}}\\[3pt]
\fbox{\begin{minipage}{0.95\linewidth}\centering Domicile \\ \footnotesize quinze éléments\end{minipage}}\\[3pt]
\fbox{\begin{minipage}{0.95\linewidth}\centering Commentaire\end{minipage}}
\end{minipage}}
\end{tabular}
\end{center}

\begin{center}
\footnotesize Agencement des blocs de la fiche patient. Reconstitution structurelle : aucun
contenu, aucune valeur.
\end{center}

\medskip

\begin{longtable}{@{}p{2.9cm}p{3.9cm}p{2.8cm}p{4.0cm}@{}}
\toprule
Identifiant & Libellé à l'écran & Type apparent & Valeur observée \\
\midrule\endfirsthead
\toprule
Identifiant & Libellé à l'écran & Type apparent & Valeur observée \\
\midrule\endhead
\multicolumn{4}{@{}l}{\textbf{Onglets}} \\
\texttt{REL-PAT.T01} & Général & onglet & — \\
\texttt{REL-PAT.T02} & Contacts & onglet & — \\
\texttt{REL-PAT.T03} & Quota soins primaires & onglet & — \\
\addlinespace
\multicolumn{4}{@{}l}{\textbf{Données principales}} \\
\texttt{REL-PAT.D01} & N° IPP & texte & — \\
\texttt{REL-PAT.D02} & Nom & texte & — \\
\texttt{REL-PAT.D03} & Nom de famille 1 & texte & — \\
\texttt{REL-PAT.D04} & Nom de famille 2 & texte & — \\
\texttt{REL-PAT.D05} & Sexe & liste & — \\
\texttt{REL-PAT.D06} & D. Nai & date & — \\
\texttt{REL-PAT.D07} & Âge & entier & — \\
\texttt{REL-PAT.D08} & Num & non\ determine & — \\
\texttt{REL-PAT.D09} & Type pièce d'identité & liste & \texttt{C} / \texttt{C.I.N} \\
\texttt{REL-PAT.D10} & N° pièce d'identité & texte & — \\
\texttt{REL-PAT.D11} & E. Civil & liste & \texttt{C} / \texttt{CÉLIBATAIRE} \\
\texttt{REL-PAT.D12} & Type patient & liste & — \\
\texttt{REL-PAT.D13} & Date photo & date & — \\
\texttt{REL-PAT.D14} & Modifié par & texte & — \\
\texttt{REL-PAT.D15} & Créé par & texte & — \\
\texttt{REL-PAT.D16} & Date d'attribution & date & — \\
\addlinespace
\multicolumn{4}{@{}l}{\textbf{Compagnie d'assurance}} \\
\texttt{REL-PAT.A01} & Compagnie d'assur. & liste & \texttt{00116} / \texttt{CNSS PIPC} \\
\texttt{REL-PAT.A02} & Police & texte & — \\
\texttt{REL-PAT.A03} & N° Assu & texte & — \\
\texttt{REL-PAT.A04} & Profession & texte & — \\
\texttt{REL-PAT.A05} & Num. inscription & texte & — \\
\texttt{REL-PAT.A06} & Date inscription & date & — \\
\addlinespace
\multicolumn{4}{@{}l}{\textbf{Domicile}} \\
\texttt{REL-PAT.H01} & Type & liste & — \\
\texttt{REL-PAT.H02} & Adresse & texte & — \\
\texttt{REL-PAT.H03} & Code postal & texte & — \\
\texttt{REL-PAT.H04} & État & liste & — \\
\texttt{REL-PAT.H05} & Ville & liste & — \\
\texttt{REL-PAT.H06} & Quartier & texte & — \\
\texttt{REL-PAT.H07} & Nationalité & liste & \texttt{504} / \texttt{MAROC} \\
\texttt{REL-PAT.H08} & Téléphone 1 & texte & — \\
\texttt{REL-PAT.H09} & Téléphone 2 & texte & — \\
\texttt{REL-PAT.H10} & Téléphone 3 & texte & — \\
\texttt{REL-PAT.H11} & Téléphone 4 & texte & — \\
\texttt{REL-PAT.H12} & Avertissements SMS & case\ a\ cocher & — \\
\texttt{REL-PAT.H13} & E-mail & texte & — \\
\texttt{REL-PAT.H14} & Avertissements e-mail & case\ a\ cocher & — \\
\texttt{REL-PAT.H15} & Environnement & liste & — \\
\addlinespace
\multicolumn{4}{@{}l}{\textbf{Né}} \\
\texttt{REL-PAT.N01} & Nom. Père & texte & — \\
\texttt{REL-PAT.N02} & Nom. Mère & texte & — \\
\texttt{REL-PAT.N03} & Lieu de naissance - État & liste & — \\
\texttt{REL-PAT.N04} & Ville & liste & — \\
\texttt{REL-PAT.N05} & Pays & liste & \texttt{504} / \texttt{MAROC} \\
\texttt{REL-PAT.N06} & Quartier & texte & — \\
\addlinespace
\multicolumn{4}{@{}l}{\textbf{Commentaire}} \\
\texttt{REL-PAT.K01} & Commentaire & zone\ texte & — \\
\addlinespace
\bottomrule
\end{longtable}

\medskip

Quatre valeurs affichées documentent la forme des nomenclatures du
système\releve{REL-PAT.D09}. Deux sont des couples code-libellé : la compagnie d'assurance et la
nationalité, cette dernière relevée deux fois — au domicile et au lieu de naissance — sous le même
couple. Deux autres sont des couples de valeurs codées : le type de pièce d'identité et l'état
civil, tous deux abrégés par une lettre suivie de son libellé développé. Le système code donc ses
références et affiche les deux faces du couple ; c'est ce qui permet, plus loin, de reconstituer
des dimensions sans inventer de libellés.

Le bloc du domicile porte quatre lignes de téléphone et deux cases d'avertissement, l'une par
message court, l'autre par courrier électronique\releve{REL-PAT.H12}. Les canaux de contact
existent donc déjà dans la fiche.

\section{Relevé — donner rendez-vous}

L'écran de prise de rendez-vous compte quarante-deux éléments, dont sept boutons de barre d'outils
hors tableau. Les trente-cinq restants se répartissent en quatre blocs : un rappel d'identité du
patient, le rendez-vous lui-même, un bloc de contrôle des modifications, et la liste d'attente des
consultations.

\begin{center}
\begin{tabular}{@{}c@{}}
\fbox{\begin{minipage}{0.86\linewidth}
\centering\small
\fbox{\begin{minipage}{0.95\linewidth}\centering Barre d'actions \\ \footnotesize sept boutons, hors tableau\end{minipage}}\\[3pt]
\fbox{\begin{minipage}{0.95\linewidth}\centering Identification du patient \\ \footnotesize dix éléments, rappelés de la fiche\end{minipage}}\\[3pt]
\fbox{\begin{minipage}{0.95\linewidth}\centering Rendez-vous \\ \footnotesize quatorze éléments\end{minipage}}\\[3pt]
\fbox{\begin{minipage}{0.46\linewidth}\centering Contrôle de modifications \\ \footnotesize huit éléments\end{minipage}}%
\hfill
\fbox{\begin{minipage}{0.46\linewidth}\centering Liste d'attente \\ \footnotesize trois colonnes\end{minipage}}
\end{minipage}}
\end{tabular}
\end{center}

\begin{center}
\footnotesize Agencement des blocs de l'écran de prise de rendez-vous. Reconstitution
structurelle : aucun contenu, aucune valeur.
\end{center}

\medskip

\begin{longtable}{@{}p{2.9cm}p{3.9cm}p{2.8cm}p{4.0cm}@{}}
\toprule
Identifiant & Libellé à l'écran & Type apparent & Valeur observée \\
\midrule\endfirsthead
\toprule
Identifiant & Libellé à l'écran & Type apparent & Valeur observée \\
\midrule\endhead
\multicolumn{4}{@{}l}{\textbf{Identification du patient}} \\
\texttt{REL-RDV.I01} & N° IPP & texte & — \\
\texttt{REL-RDV.I02} & Téléphone & texte & — \\
\texttt{REL-RDV.I03} & Compagnie & liste & — \\
\texttt{REL-RDV.I04} & N° Assu & texte & — \\
\texttt{REL-RDV.I05} & Âge & entier & — \\
\texttt{REL-RDV.I06} & Date naissance & date & — \\
\texttt{REL-RDV.I07} & C.I.N & texte & — \\
\texttt{REL-RDV.I08} & Num & non\ determine & — \\
\texttt{REL-RDV.I09} & Adresse & texte & — \\
\texttt{REL-RDV.I10} & Ville & liste & — \\
\addlinespace
\multicolumn{4}{@{}l}{\textbf{Rendez-vous}} \\
\texttt{REL-RDV.R01} & Agenda & liste & — \\
\texttt{REL-RDV.R02} & Activité & liste & — \\
\texttt{REL-RDV.R03} & Origine & liste & \texttt{AU} / \texttt{AUTRES HÔPITAUX} \\
\texttt{REL-RDV.R04} & Hôpital/C.S. & liste & — \\
\texttt{REL-RDV.R05} & Médecin ext. & texte & — \\
\texttt{REL-RDV.R06} & Service ext. & texte & — \\
\texttt{REL-RDV.R07} & Observations & zone\ texte & — \\
\texttt{REL-RDV.R08} & Date rendez-vous & horodatage & \texttt{07/28/2026 10:57:07 AM} \\
\texttt{REL-RDV.R09} & Rendez-vous supplémentaire & case\ a\ cocher & — \\
\texttt{REL-RDV.R10} & Type d'attention & liste & — \\
\texttt{REL-RDV.R11} & État & liste & \texttt{En instance} \\
\texttt{REL-RDV.R12} & Durée & entier & \texttt{0} \\
\texttt{REL-RDV.R13} & Date réception & horodatage & — \\
\texttt{REL-RDV.R14} & Imprimer données & case\ a\ cocher & — \\
\addlinespace
\multicolumn{4}{@{}l}{\textbf{Contrôle de modifications}} \\
\texttt{REL-RDV.C01} & Créé par & texte & — \\
\texttt{REL-RDV.C02} & Date création & horodatage & \texttt{07/28/2026 10:57:06 AM} \\
\texttt{REL-RDV.C03} & Modifié par & texte & — \\
\texttt{REL-RDV.C04} & Date mod. & horodatage & — \\
\texttt{REL-RDV.C05} & Confirmé par & texte & — \\
\texttt{REL-RDV.C06} & Date conf. & horodatage & — \\
\texttt{REL-RDV.C07} & Annulé par & texte & — \\
\texttt{REL-RDV.C08} & Date annul. & horodatage & — \\
\addlinespace
\multicolumn{4}{@{}l}{\textbf{Liste d'attente des consultations}} \\
\texttt{REL-RDV.L01} & Service & colonne & — \\
\texttt{REL-RDV.L02} & Agenda & colonne & — \\
\texttt{REL-RDV.L03} & Activité & colonne & — \\
\addlinespace
\bottomrule
\end{longtable}

\medskip

Le bloc du rendez-vous porte l'origine de la demande sous forme de couple code-libellé, et la
valeur observée désigne un adressage par une autre structure
hospitalière\releve{REL-RDV.R03}. L'état du rendez-vous à la création et sa durée sont également
lisibles\releve{REL-RDV.R11}. Le bloc de contrôle des modifications porte quatre couples
d'auteur et d'horodatage — création, modification, confirmation, annulation — dont la portée est
traitée à la section suivante.

Les dix éléments du bloc d'identification ne sont pas des champs de saisie propres à cet écran :
ils rappellent des valeurs de la fiche patient. Le modèle de données ne les enregistre donc pas
deux fois.

\section{Trois observations exploitables}

Trois relevés de cet écran dépassent la description : ils fondent une grandeur, une distinction et
un diagnostic que le reste du rapport reprend. Chacun est une observation, et rien d'autre.

\textbf{L'écart entre la création et le rendez-vous est le délai
d'obtention}\releve{REL-RDV.C02}. Le bloc de contrôle des modifications horodate la création ; le
bloc du rendez-vous horodate la date convenue. Leur différence est le délai qu'un patient attend,
grandeur que ni l'un ni l'autre des deux champs ne porte seul. Sur le rendez-vous observé, cette
différence vaut une seconde : la création à \texttt{10:57:06} et le rendez-vous à
\texttt{10:57:07} du même jour, soit un rendez-vous donné pour l'instant même. Le cas est extrême
et il est instructif : il montre que le système accepte un délai nul, donc que la grandeur est
bornée à zéro et non strictement positive. Le chapitre~\ref{chap:activite} en fait un indicateur.

\textbf{Un couple de champs sépare l'annulation de l'absence}\releve{REL-RDV.C07}. Le bloc de
contrôle porte un champ d'auteur de confirmation et un champ d'auteur d'annulation, chacun avec
son horodatage. Sans ce couple, un rendez-vous non honoré serait un état unique. Avec lui, deux
situations se distinguent : l'annulation déclarée, qui laisse au service le temps de réaffecter le
créneau, et l'absence non prévenue, qui le perd. La distinction n'est pas cosmétique — elle sépare
deux taux dont les remèdes diffèrent, et le chapitre~\ref{chap:activite} les mesure séparément.

\textbf{La recherche pondérée est un aveu}\releve{REL-IPP.O03}. Un système qui offre de chercher
un patient par similarité, et qui affiche une probabilité en regard de chaque résultat, admet que
la recherche exacte ne suffit pas. Elle ne suffit pas parce qu'une même personne y possède
plusieurs dossiers. La fonction existe dans le produit installé ; ce que le relevé ne dit pas,
c'est si le service l'emploie. Le chapitre~\ref{chap:qualite} construit un rapprochement
probabiliste sur ce constat, et le confronte à une collision exacte pour mesurer ce que la
similarité apporte.

Une quatrième observation, mineure et têtue, court à travers les trois précédentes : les
horodatages s'affichent au format mois-jour-année avec un indicateur de méridien, dans une
interface par ailleurs française\releve{REL-RDV.R08}. La date observée le montre sans ambiguïté —
le quantième vaut vingt-huit, ce qui exclut toute lecture jour-mois, et lève l'incertitude qu'une
date antérieure au treizième jour du mois aurait laissée entière.

Ce détail d'affichage n'est pas cosmétique. Une extraction du système rend ses dates dans la forme
où l'écran les compose ; toute date entrant dans la chaîne de données doit donc être convertie
depuis ce format, et la conversion doit être écrite une fois et vérifiée, faute de quoi un jour et
un mois s'échangent silencieusement onze fois sur douze sans qu'aucune valeur ne devienne
invalide. C'est une contrainte de chargement, reprise au chapitre~\ref{chap:architecture}.

Ces quatre observations partagent une propriété qu'il faut nommer, parce qu'elle gouverne la
suite. Aucune ne résulte d'un entretien ni d'une documentation : chacune est lisible sur un écran,
et chacune est consignée sous un identifiant que le lecteur peut retrouver dans les tableaux des
sections précédentes. Elles ne disent rien de la fréquence des situations qu'elles révèlent — un
seul rendez-vous a été observé, et un cas n'est pas une distribution. Ce qu'elles établissent est
qu'une grandeur est calculable, qu'une distinction est enregistrée, qu'une fonction existe et
qu'un format s'impose. C'est précisément ce dont la modélisation a besoin : non pas des valeurs,
mais la certitude que les champs qui les porteraient existent.

% Ce chapitre repose sur :
%   DOC: s30
%   OBS: presentation-du-besoin, rattachement-du-travail, organisation-en-phases
%   HYP: (aucun)
%   SECTIONS: 6

\chapter{Cadrage et méthodologie}\label{chap:cadrage}

\section{Le besoin exprimé}

Le travail répond à un cahier de charges portant sur l'activité du service d'accueil et
d'admission dans son ensemble, et non sur le seul processus observable depuis un poste. Ce
périmètre est déterminant pour la suite : il oblige à documenter quatre processus auxquels
l'observation directe n'a pas donné accès, et il explique la place que tient dans ce projet la
reconstruction par voie documentaire.

Le document lui-même ne figure pas au dépôt et n'a pas vocation à y figurer. Ce qui y figure, en
revanche, ce sont trois de ses exigences, consignées à l'endroit où elles ont été mises en défaut :
tout indicateur affiché doit être recalculé depuis les tables de faits, le tableau de bord ne doit
lire que la base, et les valeurs affichées doivent être lisibles par un lecteur qui ne connaît pas
les codes du système d'origine.

Les trois sont tenues en partie et manquées en partie, et les écarts sont écrits plutôt que
contournés. Six indicateurs ne sont pas recalculés depuis les faits, dont cinq parce qu'aucune table
de faits ne porte la matière — il n'existe ni fait des créances ni fait des relances — et un parce
que le recalcul changerait sa valeur en changeant sa population. La capacité litière, dont dépendent
quatre indicateurs de séjour, n'est lisible nulle part dans la base et transite par une table de
paramètres portant sa provenance. Enfin quatre dimensions sur six affichent leur code nu, faute de
libellé qu'une source établisse — et l'abstention est délibérée : un code nu est illisible et le
montre, un libellé inventé est lisible et le cache.

\aRediger{presentation-du-besoin}{la façon dont le besoin a été présenté au début du stage, et ce
que l'encadrement a précisé oralement sur les attentes et les priorités}\releve{presentation-du-besoin}

\section{Périmètre d'observation}

Le système d'information cloisonne l'accès par profil applicatif, et le service n'accorde à un
stagiaire qu'un seul des cinq profils que porte l'écran de démarrage. \textbf{Un processus a donc été
observé, quatre ont été reconstruits par voie documentaire.} Ce n'est pas une commodité : c'est la
contrainte à partir de laquelle tout le dispositif de provenance a été bâti.

Une observation et une reconstruction n'ont pas la même valeur, et les confondre reviendrait à
présenter comme vu ce qui a été déduit. Le projet ne les confond pas : \textbf{chaque colonne du modèle de
données porte une étiquette de provenance, et l'exigence de preuve diffère par étiquette.}

\begin{center}
\begin{tabular}{@{}p{1.6cm}p{5.0cm}p{1.6cm}p{5.5cm}@{}}
\toprule
Étiquette & Sens & Entrées & Ce que la preuve doit être \\
\midrule
\texttt{OBS} & champ observé à l'écran & \textbf{81} & une référence de relevé, et le libellé reproduit à l'identique \\
\addlinespace
\texttt{DOC} & champ déduit d'une source & \textbf{72} & un identifiant du registre des sources \\
\addlinespace
\texttt{HYP} & champ posé sans preuve externe & \textbf{22} & la mention explicite qu'aucune preuve n'existe \\
\bottomrule
\end{tabular}
\end{center}

\begin{center}
\footnotesize Les trois étiquettes de provenance et leur exigence de preuve, sur les 175 colonnes de
la couche source.
\end{center}

\medskip

Les vingt-deux colonnes hypothétiques portent toutes la mention d'absence de preuve externe, et ce
décompte est mesuré, non supposé. La proportion elle-même est un résultat : \textbf{un peu moins de la
moitié des colonnes est observée}, et le lecteur sait laquelle.

Le même dispositif s'étend à la prose du présent rapport. Chaque affirmation y relève d'une des trois
mêmes catégories — une source documentaire, un relevé d'observation, une convention posée —, chaque
fichier de chapitre déclare en tête ce sur quoi il repose, et un contrôle vérifie que la déclaration
coïncide exactement avec ce que le fichier porte, dans les deux sens.

\textbf{Aucune image du système d'information n'entre dans ce rapport, et le motif tient en une
phrase :} les écrans d'un dossier patient en production affichent des identifiants de session, des
libellés de compte et la structure des dossiers d'un établissement réel, et le dépôt de ce projet est
public. Ce que les images auraient montré, le relevé de champs le porte sous une forme plus
exploitable — quatre écrans, cent quarante-trois éléments, chacun avec son identifiant, son libellé
exact, son type apparent et sa valeur observée le cas échéant. Deux schémas d'écran redessinés
restituent l'agencement des blocs, sans aucun contenu. La règle est tenue par deux dispositifs
indépendants : une exclusion dans la configuration du gestionnaire de versions, et un contrôle
bloquant qui échoue si un fichier suivi porte une extension d'image sous le répertoire du rapport ou
à la racine — le premier seul ne protégerait pas d'une indexation forcée.

\section{La contrainte de confidentialité et la décision de générer}

Une décision de cadrage domine tout le projet : \textbf{aucune donnée réelle ne sort du service.} Elle
n'est pas négociable et elle n'a pas été négociée. Elle interdit d'extraire un échantillon, fût-il
anonymisé, et elle interdit donc de bâtir une chaîne de données sur des lignes venues du système
observé.

La conséquence est qu'un jeu de données est produit. Sa volumétrie n'est pas inventée : elle est
relevée par établissement nommé dans le recueil statistique national \cite{s30} et mise à l'échelle
de la période. Ses relations internes ne sont pas davantage improvisées — elles sont injectées
délibérément, une par une, et chacune est consignée avec son paramètre, son statut de provenance et
sa conséquence sur ce que le tableau de bord donne à voir.

\textbf{Ce que la génération permet.} Elle permet de démontrer une chaîne complète sur des volumes
réalistes, d'y injecter des défauts dont on connaît la vérité terrain, et donc de mesurer ce qu'un
traitement retrouve. Un rapprochement probabiliste évalué contre des doublons dont on sait
lesquels sont vrais est mesurable ; le même rapprochement sur des données réelles ne le serait pas.

\textbf{Ce qu'elle interdit de conclure, et c'est le point le plus important de cette section.} Aucune
grandeur produite par la chaîne n'est une mesure de l'établissement. Une relation injectée se
retrouve nécessairement à la sortie : la lire comme un résultat serait circulaire. Le projet écrit
cette limite là où elle se lit — le registre des relations injectées porte, pour chacune, ce que la
page affiche et ce qu'on n'a pas le droit d'en conclure. Le chapitre de recommandations en tire une
conséquence de forme, en séparant systématiquement ce qui établit un constat de ce qu'un indicateur
démontre.

\section{Démarche}

Deux principes gouvernent le travail, et l'un et l'autre se lisent dans le dépôt plutôt que dans une
déclaration d'intention.

\textbf{Mesurer avant de corriger.} Aucune correction n'est proposée avant que sa cause soit
mesurée, et aucun chiffre n'est écrit sans la commande qui le produit. Le principe se prolonge
au-delà de l'écriture : les enregistrements de décision portent, pour plusieurs d'entre eux, une
prémisse explicitement datée et la consigne de la re-mesurer avant toute modification qui la
mettrait en jeu. Une décision n'y est jamais présentée comme définitive, mais comme vraie sous des
conditions nommées.

\textbf{Vérifier par mutation.} Une propriété nouvellement affirmée n'est pas tenue pour vérifiée
parce qu'un contrôle est vert : le contrôle est délibérément mis en défaut, on s'assure qu'il
rougit, et l'état initial est restauré. Un contrôle vert qui reste vert quand on casse ce qu'il
surveille ne surveille rien.

La méthode a payé, et il faut dire ce qu'elle a trouvé plutôt que d'en vanter le principe. Elle a
révélé qu'un contrôle de correspondance passait par accident sur un fichier dépourvu de déclaration
— ensemble déclaré vide contre ensemble porté vide —, ce qui a fait ajouter l'exigence d'une
déclaration explicite. Elle a révélé ensuite que le retrait d'une section entière à l'intérieur d'un
chapitre n'était détecté que si cette section emportait une source qu'aucune autre ne portait :
mesuré sur huit sections, six étaient détectées et deux ne l'étaient pas, ce qui a fait déclarer le
décompte de sections. Elle a enfin établi qu'une vérification prévue ne pouvait pas mordre — le
convertisseur bibliographique ne lisant jamais le champ que la mutation visait —, cas où la
vérification est remplacée par celle qui a un sens plutôt que forcée.

\aRediger{rattachement-du-travail}{le rattachement du travail effectué aux phases du projet, et ce
que chaque phase a produit}\releve{rattachement-du-travail}

\section{Outils et justification}

Chaque brique de la pile technique est retenue pour un motif écrit, et non par habitude.

\begin{center}
\begin{tabular}{@{}p{3.4cm}p{4.0cm}p{6.4cm}@{}}
\toprule
Brique & Rôle & Motif du choix \\
\midrule
Serveur de base de données relationnelle & couche de persistance & un moteur client-serveur, non un moteur embarqué : trois services attaquent la même base par le réseau, avec des identifiants distincts, et le chargement écrit pendant que le tableau de bord lit \\
\addlinespace
Schéma en étoile & couche analytique & des faits entourés de dimensions, forme que l'interrogation analytique exploite et que le dossier doit démontrer \\
\addlinespace
Outil de transformation déclaratif & modélisation & matérialisation en vues, configuration de connexion hors dépôt, macros de conversion adossées aux formats mesurés à l'extraction \\
\addlinespace
Bibliothèque de rapprochement probabiliste & identités & un modèle de vraisemblance dont les paramètres s'estiment sur les données, plutôt qu'une collision exacte qui ne retrouve que l'identique \\
\addlinespace
Moteur en mémoire pour le rapprochement & exécution & le rapprochement écrit des tables intermédiaires nombreuses ; les tenir hors du serveur évite d'y laisser des résidus \\
\addlinespace
Orchestrateur de graphes de tâches & ordonnancement & le conteneur ordonnance sans exécuter : les tâches sont des commandes qui s'exécutent dans l'environnement du projet \\
\addlinespace
Bibliothèque d'application web analytique & tableau de bord & un service qui lit un schéma dédié, rafraîchi par échange de noms en une transaction \\
\bottomrule
\end{tabular}
\end{center}

\begin{center}
\footnotesize Pile technique et motif de chaque choix. Chacun est consigné dans un enregistrement
de décision distinct, avec la mesure qui l'a tranché et ce qui l'aurait invalidé.
\end{center}

\medskip

Deux traits de cette liste méritent d'être relevés. Le premier est que le motif d'un choix n'est
presque jamais la performance : c'est la propriété que le dossier doit démontrer, ou la contrainte
qu'il faut ne pas violer. Le second est que chaque enregistrement porte une rubrique disant ce qui
aurait invalidé la décision — un choix qu'on ne sait pas décrire comme réfutable n'est pas un choix,
c'est une préférence.

\section{Organisation du projet en phases}

\aRediger{organisation-en-phases}{les phases du projet, nommées par ce qu'elles font, avec pour
chacune son objectif et ses livrables, suivies d'une figure restituant leurs dépendances — aucun
calendrier, aucun relevé de jours}\releve{organisation-en-phases}

\chapter{Conception du jeu de données}

\section{Pourquoi un jeu simulé}

\section{La volumétrie relevée et sa mise à l'échelle}

\section{Les relations injectées et leur circularité}

\section{Les défauts injectés}

% Ce chapitre repose sur :
%   DOC: (aucun)
%   OBS: (aucun)
%   HYP: (aucun)
%   SECTIONS: 4

\chapter{Architecture de la chaîne de données}\label{chap:architecture}

\section{Les couches et leur rôle}

\section{Le chargement et la quarantaine}

\section{La modélisation dimensionnelle}

\section{L'orchestration}

\chapter{Qualité des données et rapprochement d'identités}

\section{Les contrôles de qualité et leurs seuils}

\section{Le rapprochement probabiliste}

\section{L'évaluation contre la vérité terrain}

\chapter{Analyse de l'activité}

\section{Les indicateurs de séjour}

\section{Les rendez-vous et le délai d'obtention}

\section{Les urgences}

\section{La facturation et le recouvrement}

% Ce chapitre repose sur :
%   DOC: (aucun)
%   OBS: (aucun)
%   HYP: (aucun)
%   CHI: (aucun)
%   SER: (aucune)
%   SECTIONS: 3

\chapter{Le tableau de bord}

\section{La composition des pages}

\section{Le registre des indicateurs}

\section{Ce que l'écran déclare de ses propres limites}

% Ce chapitre repose sur :
%   DOC: s12, s13, s20, s24, s30
%   OBS: REL-PAT.H12, REL-PAT.H14, REL-IPP.O03, REL-IPP.R12
%   HYP: absentéisme-par-spécialité, complétude-par-champ, ancienneté-des-créances,
%   HYP: répartition-des-niveaux-de-tri
%   CHI: (aucun)
%   SECTIONS: 5

\chapter{Recommandations}\label{chap:recommandations}

Ce chapitre propose cinq actions. Chacune tient en un bloc de forme fixe, et ce bloc sépare deux
choses que l'on confond ordinairement.

\textbf{Ce qui établit le constat} est ce qui autorise à dire qu'un problème existe. Sur un jeu de
données produit pour ce projet, cela ne peut pas être une valeur affichée par le tableau de bord :
les relations qui gouvernent ces valeurs y ont été injectées, et les relire à la sortie serait
circulaire. Le constat vient donc d'une source extérieure au jeu — un rapport de contrôle, une
publication statistique, une observation au poste — ou bien il n'existe pas, et le rapport le dit.

\textbf{Ce que l'indicateur démontre} est autre chose, et ce n'est pas rien : il démontre que la
chaîne de données sait produire le nombre, à la maille utile, à partir des tables. Le jour où elle
sera branchée sur les données réelles du service, le nombre sera un constat. Aujourd'hui, il est une
capacité.

\textbf{Aucun effort n'est chiffré et aucun effet n'est quantifié.} Un effort en jours et un gain en
points seraient des nombres sans mesure qui les produise, et ce rapport ne s'en autorise pas. Chaque
bloc dit ce que l'action suppose — quel champ existe déjà, quel processus change — et par quelle
mesure on constaterait son effet. L'effet lui-même, lorsqu'il est énoncé, l'est comme une attente.

\section{Sur les rendez-vous}

\textbf{Le constat.} Le délai d'obtention d'un rendez-vous varie fortement d'une spécialité à
l'autre, et une part des créneaux est perdue par des patients qui ne se présentent pas. Le rapport
de la Cour des comptes consacré au centre hospitalier établit l'un et l'autre pour cet établissement
même : il relève des délais moyens de rendez-vous par spécialité, et l'absence de tout état
comparatif entre rendez-vous demandés et consultations réalisées \cite{s20}. L'ordre de grandeur de
l'absentéisme ambulatoire est documenté par ailleurs — une moyenne mondiale de 23,5~\% et une
fourchette de 23 à 33~\% selon les contextes \cite{s24}.

\textbf{Ce que l'indicateur démontre.} Le tableau de bord calcule le délai d'obtention en médiane et
au quatre-vingt-dixième centile par activité, le taux d'absence et le taux d'annulation par
activité. Il démontre que ces quatre grandeurs se produisent depuis le fait des rendez-vous, à la
maille où une décision se prend. Leurs valeurs, en revanche, sont fixées par des paramètres
injectés\convention{absentéisme-par-spécialité} : le classement des spécialités par absentéisme est
imposé, et aucune conclusion sur la force du lien entre délai et absence ne peut en être tirée.

\textbf{L'action.} Rappeler les patients dont le rendez-vous approche, en commençant par les
activités où le délai est le plus long, et ajuster le nombre de créneaux ouverts au taux d'absence
mesuré sur chaque activité une fois la chaîne branchée sur les données réelles.

\textbf{Ce qu'elle suppose.} Le canal existe déjà : la fiche patient porte quatre lignes de
téléphone, une adresse électronique, et deux cases d'avertissement — l'une par message court,
l'autre par courrier électronique\releve{REL-PAT.H12}. Le relevé note que ces cases existent et ne
sont pas exploitées par le service\releve{REL-PAT.H14}. L'action ne demande donc aucun développement
logiciel : elle demande que le champ soit renseigné à l'admission et qu'un envoi soit déclenché.

\textbf{L'effet attendu et sa mesure.} On attend une baisse du taux d'absence sur les activités
ciblées, et une réduction du délai médian sur ces mêmes activités par récupération des créneaux
perdus. La mesure existe : le taux d'absence et le délai d'obtention par activité, comparés avant et
après, à activité constante. La distinction entre absence et annulation, que le système enregistre
par un couple de champs distinct, permet de vérifier que la baisse vient bien des absences et non
d'un report de qualification.

\section{Sur les urgences}

\textbf{Le constat.} Une part importante des passages aux urgences relève d'une consultation
ordinaire. L'avis du Conseil Économique, Social et Environnemental l'établit à l'échelle nationale :
sur les 6~482~185 consultations et soins d'urgence assurés chaque année par le secteur public,
\textbf{64~\% sont des consultations non urgentes} et 10~\% seulement des urgences vitales
\cite{s12}. Une étude conduite dans un hôpital provincial marocain mesure 30,7~\% de consultations
non appropriées sur 410 patients \cite{s13}. Les deux sources concordent sur l'existence du
phénomène ; elles divergent sur son ampleur, l'une comptant les motifs de recours et l'autre le
caractère approprié de la venue.

\textbf{Ce que l'indicateur démontre.} Le tableau de bord calcule la part des passages classés aux
deux niveaux de tri les moins graves, le délai de prise en charge médicale par niveau, la durée de
passage et l'orientation de sortie. Il démontre que la file d'attente se décrit par niveau de
gravité et que le délai s'y mesure. La pyramide de gravité affichée est en revanche
imposée\convention{répartition-des-niveaux-de-tri}, y compris la part relevant d'une consultation
ordinaire : c'est un paramètre, non une observation.

\textbf{L'action.} Ouvrir une filière courte pour les deux niveaux de tri les moins graves, prise en
charge séparément de la file principale, et réorienter en amont vers les établissements de soins de
santé primaires les venues qui relèvent d'eux.

\textbf{Ce qu'elle suppose.} Un tri formalisé à l'accueil, dont le résultat est enregistré : sans
niveau de tri saisi, ni la filière ni sa mesure n'existent. Elle suppose aussi une capacité d'accueil
en amont, que ce rapport ne mesure pas.

\textbf{L'effet attendu et sa mesure.} On attend une baisse du délai de prise en charge des niveaux
les plus graves, la file courte cessant de les précéder. La mesure est le délai de prise en charge
médicale par niveau de tri, en médiane et au quatre-vingt-dixième centile : c'est sur les niveaux
graves qu'il faut la lire, et non sur la moyenne d'ensemble, qu'une filière courte améliore
mécaniquement.

\section{Sur la facturation et le recouvrement}

\textbf{Le constat.} Il est établi pour cet établissement, et par une source extérieure au jeu de
données. La Cour des comptes relève \textbf{50~876~365 dirhams de créances non recouvrées} au
31 décembre 2016, \textbf{57~\% des dossiers non facturés} en 2015, et \textbf{3~477 consultations
payées sur 160~659} aux urgences en 2016 \cite{s20}. Deux problèmes distincts s'y lisent : des
créances qui vieillissent sans être recouvrées, et des épisodes de soin qui ne donnent lieu à aucune
facture.

\textbf{Ce que l'indicateur démontre.} Le tableau de bord calcule le taux de recouvrement des
créances, leur répartition par tranche d'ancienneté, l'aboutissement des relances, le taux
d'encaissement, et le nombre d'épisodes sans facture rattachée par famille et par mois. Il démontre
que la chaîne sait rapprocher un épisode de sa facture et une facture de son encaissement — ce qui
est précisément ce dont l'absence était relevée. Les montants et l'ancienneté affichés sont
imposés\convention{ancienneté-des-créances}, et la priorisation qu'on en tirerait serait circulaire.

\textbf{L'action.} Prioriser les relances par ancienneté décroissante et par type de débiteur, et
produire chaque mois la liste des épisodes sans facture rattachée pour reprise par le service.

\textbf{Ce qu'elle suppose.} Que l'épisode et la facture portent une clé commune exploitable, ce que
le modèle de données démontre possible. Elle suppose surtout un changement de processus et non
d'outil : la créance naît au service d'accueil, où la facture est établie et le règlement encaissé à
la sortie, tandis que le recouvrement relève du pôle administratif — la liste doit donc traverser
cette frontière pour servir.

\textbf{L'effet attendu et sa mesure.} On attend une hausse du taux de recouvrement et une baisse de
la part des créances les plus anciennes. La mesure du rendement de la relance existe — nombre de
relances émises, nombre ayant abouti, et leur rapport — et elle doit être lue avant d'augmenter le
volume des relances plutôt qu'après.

\section{Sur la qualité du fichier patient}

\textbf{Le constat, et il commence par ce qu'on ne sait pas.} Le système n'expose pas quels champs
de la fiche patient sont obligatoires à la saisie et quels champs sont facultatifs. L'observation au
poste ne l'a pas relevé, et le relevé de champs le déclare explicitement : le caractère obligatoire
ou facultatif n'a été relevé sur aucun écran, et n'est donc affirmé nulle part. Recommander de rendre
certains champs obligatoires supposerait de savoir lesquels ne le sont pas. \textbf{C'est le premier
constat, et il fonde la première action.}

Le second constat est établi, lui, et par le système lui-même : une même personne y possède
plusieurs dossiers. L'écran de recherche d'identifiant porte un onglet de recherche pondérée sur
l'index maître des patients\releve{REL-IPP.O03}, et ses résultats s'accompagnent d'une colonne de
probabilité\releve{REL-IPP.R12}. Un système qui offre de chercher un patient par similarité admet
que la recherche exacte ne suffit pas. Le rapport de la Cour des comptes le confirme de l'extérieur,
relevant l'attribution de plusieurs identifiants au même patient et recommandant l'enregistrement de
la pièce d'identité \cite{s20}.

\textbf{Ce que l'indicateur démontre.} Le tableau de bord calcule le taux de renseignement de chaque
colonne de chaque table, le taux de mise en quarantaine, le nombre de groupes de patients partageant
exactement les mêmes valeurs d'identité, et le nombre de grappes formées par le rapprochement
probabiliste. Il démontre que la complétude se mesure colonne par colonne et que l'ampleur du
doublon se chiffre. Le classement des champs par complétude est en revanche
imposé\convention{complétude-par-champ}, et — contrairement aux trois recommandations précédentes —
\textbf{aucune source extérieure ne vient l'établir}. C'est la recommandation la moins appuyée du
chapitre, et il vaut mieux l'écrire que le taire.

\textbf{L'action, en deux temps.} D'abord établir quels champs sont facultatifs, par relevé du
paramétrage du système ou à défaut par mesure de la complétude réelle sur un extrait ; puis rendre
obligatoires ceux qui portent l'identité et le contact. Ensuite, déplacer le contrôle de doublon :
le rapprochement doit s'exécuter au moment où l'identifiant est créé, et non après coup sur un
fichier déjà pollué.

\textbf{Ce qu'elle suppose.} Pour le second temps, presque rien : la fonction existe déjà dans le
produit installé, sous la forme de la recherche pondérée relevée à l'écran. Ce qui manque n'est pas
l'outil mais son insertion dans le geste d'accueil — la recherche doit être obligatoire avant
création, non facultative après.

\textbf{L'effet attendu et sa mesure.} On attend une baisse du nombre de groupes d'identité formés
par collision exacte, et une hausse du taux de renseignement des champs rendus obligatoires. Les deux
mesures existent, et la seconde vérifie la première : un champ rendu obligatoire mais renseigné par
une valeur de remplissage se verrait à la complétude sans améliorer le rapprochement.

\section{Sur les indicateurs que le système ne permet pas de produire}

\textbf{Cette recommandation ne dépend d'aucune donnée générée}, et c'est la seule du chapitre dans
ce cas. Elle porte sur des champs absents, fait qui ne se lit pas dans des valeurs mais dans une
confrontation entre ce que la statistique nationale attend d'un hôpital nommé et ce que la chaîne
peut calculer.

\textbf{Le constat.} Le recueil statistique national publie, pour chaque établissement nommé, une
liste précise d'indicateurs \cite{s30}. Trois d'entre eux ne sont pas calculables, et pour chacun le
champ manquant se nomme.

\begin{center}
\begin{tabular}{@{}p{4.4cm}p{4.6cm}p{5.0cm}@{}}
\toprule
Indicateur attendu & Champ manquant & Ce qu'il débloquerait \\
\midrule
Capacité litière existante & une table de structure de l'établissement, portant les lits installés & la distinction entre lits installés et lits en service, et le suivi de son évolution \\
\addlinespace
Consultations par médecin & un référentiel des médecins affectés & le rapport de l'activité à l'effectif, seul indicateur de productivité que la nomenclature attende \\
\addlinespace
Accouchements & un marqueur d'accouchement sur le séjour & le décompte des accouchements, et le taux de césarienne s'il devenait pertinent \\
\bottomrule
\end{tabular}
\end{center}

\begin{center}
\footnotesize Les trois indicateurs attendus par la statistique nationale que la chaîne ne peut pas
produire, et le champ dont l'absence l'en empêche.
\end{center}

\medskip

\textbf{Une distinction gouverne ce tableau et doit être respectée.} Un indicateur non calculable
faute de champ est un défaut du système d'information, et il appelle une recommandation. Un
indicateur sans objet pour ce site est une propriété de l'établissement, et il n'appelle aucun
reproche au logiciel : les interventions chirurgicales et le taux de césarienne relèvent de cette
seconde catégorie, l'établissement n'exerçant pas d'activité chirurgicale. Confondre les deux
reprocherait au produit une absence de champ là où c'est l'hôpital qui n'a pas l'activité.

\textbf{L'action.} Ajouter les trois champs, en commençant par le référentiel des médecins, dont
dépend le seul indicateur de productivité attendu.

\textbf{Ce qu'elle suppose.} Une intervention sur le paramétrage du produit ou une demande à
l'éditeur, selon que les champs existent et ne sont pas alimentés ou qu'ils n'existent pas — ce que
l'observation conduite ici ne permet pas de trancher.

\textbf{L'effet attendu et sa mesure.} On attend que les trois indicateurs deviennent calculables,
et l'effet se constate sans ambiguïté : la valeur existe ou elle n'existe pas. La confrontation entre
la nomenclature attendue et ce que la chaîne produit, déjà conduite une fois, se refait à l'identique
et donne la réponse.


% Ce chapitre repose sur :
%   DOC: (aucun)
%   OBS: apport-personnel
%   HYP: (aucun)
%   CHI: conclusions-avec-relation, conclusions-sans-relation, exigences-calculables, exigences-champs-manquants, exigences-non-calculables, exigences-rubriques
%   CHI: exigences-sans-objet, fichiers-de-controle, limites-des-faits, modeles-total, relations-injectees, relations-non-reprises
%   CHI: releve-champs-non-employes, releve-champs-total, taches-graphe, tdb-indicateurs, tdb-pages
%   SER: (aucune)
%   SECTIONS: 0

\chapter*{Conclusion générale}
\addcontentsline{toc}{chapter}{Conclusion générale}
% Un chapitre étoilé ne met pas à jour la marque de l'en-tête courant : sans cette ligne,
% les pages de ce texte portent l'en-tête de ce qui le précède — mesuré, « TABLE DES
% MATIÈRES » sur la deuxième page de l'introduction et « CHAPITRE 9 » sur celle de la
% conclusion. Le défaut ne se voit qu'en regardant le document composé.
\markboth{Conclusion générale}{Conclusion générale}

\section*{Ce qui a été fait}

Une chaîne de données complète a été construite, du fichier déposé jusqu'à l'écran :
\chiffre{modeles-total} modèles de transformation, \chiffre{taches-graphe} tâches orchestrées, un
rapprochement probabiliste d'identités évalué contre une vérité terrain connue, et un tableau de
bord de \chiffre{tdb-pages} pages portant \chiffre{tdb-indicateurs} indicateurs. Elle se rejoue en
entier, et \chiffre{fichiers-de-controle} fichiers de contrôle la surveillent.

Ce n'est pas le volume de ce travail qui en fait l'intérêt, mais \textbf{ce qu'il refuse
d'affirmer}. Chaque nombre du rapport vient d'une commande qui le produit ; chaque affirmation porte
l'étiquette de ce qui l'étaye — un document, une observation, ou une convention posée ; chaque
relation injectée est déclarée, et le rapport dit laquelle des conclusions la reproduit.
\textbf{Le résultat le plus directement utilisable est la correspondance avec ce que la
nomenclature nationale attend} : sur \chiffre{exigences-rubriques} rubriques attendues d'un
établissement nommé, \chiffre{exigences-calculables} sont calculables par la chaîne.

\section*{Ce qui n'a pas été fait, et pourquoi}

\textbf{Des rubriques de la nomenclature restent hors d'atteinte, et aucune ne l'est par manque de
temps.} Elles sont \chiffre{exigences-non-calculables}, bloquées chacune par un champ absent du
système d'information, et \chiffre{exigences-champs-manquants} champs suffiraient à les débloquer.

Il s'y ajoute \chiffre{exigences-sans-objet} rubriques sans objet, et la distinction compte : elles
ne manquent pas au logiciel, c'est l'établissement qui n'exerce pas l'activité.

\textbf{Le modèle dimensionnel porte \chiffre{limites-des-faits} limites documentées plutôt que
contournées}, et toutes tiennent à la même cause : la source ne porte pas ce qu'il faudrait pour
aller plus loin. Fabriquer la distinction manquante aurait produit une donnée qu'aucune ligne
n'établit.

\textbf{Sur \chiffre{relations-injectees} relations injectées,
\chiffre{relations-non-reprises} ne sont reprises par aucune conclusion}, le plus souvent parce que
la colonne qui les porterait n'est calculée dans aucune couche aval. Le régime de couverture, l'âge
et l'adressage existent dans la source et s'arrêtent en chemin.

\textbf{Le relevé d'observation, lui, est employé en entier ou presque.} Sur
\chiffre{releve-champs-total} éléments relevés au poste,
\chiffre{releve-champs-non-employes} seulement ne sont cités nulle part — et ce sont tous des
boutons. Aucun champ de donnée observé n'est resté inemployé.

\textbf{Enfin, aucune conclusion du chapitre d'analyse n'établit un fait sur le service.} Sur les
vingt-deux qu'il porte, \chiffre{conclusions-avec-relation} reproduisent une relation injectée et
\chiffre{conclusions-sans-relation} ne viennent d'aucune. Ce n'est pas un défaut d'exécution :
c'est la conséquence nécessaire d'un jeu produit.

\section*{Les pistes qui n'ont pas été retenues}

Trois pistes ont été écartées, et les trois pour le même genre de raison. \textbf{Inventer les
libellés manquants} : un code nu est illisible et le montre, un libellé inventé est lisible et cache
qu'il est inventé. \textbf{Régénérer le jeu de données} pour effacer le désaccord des deux faits sur
la fin d'un épisode : un jeu qu'on régénère jusqu'à ce qu'il soit propre n'apprend plus rien.
\textbf{Illustrer le tableau de bord par des captures d'écran} : le dépôt interdit toute image du
système d'information.

\section*{Ce qu'il faudrait pour passer à des données réelles}

La chaîne est écrite pour lire des fichiers, non pour lire ce générateur : la question n'est pas
de savoir si elle fonctionnerait, mais \textbf{ce qui changerait de nature}. Trois choses.

\textbf{La vérité terrain disparaît.} Précision et rappel se calculent aujourd'hui parce qu'on sait
lesquels des doublons sont vrais ; sur des données réelles, cette connaissance n'existe pas.
L'évaluation changerait donc de nature — il faudrait un échantillon annoté à la main, avec son coût
et son incertitude propres — et les grandeurs publiées au chapitre du rapprochement cesseraient
d'être mesurables telles quelles.

\textbf{Les distributions ne sont plus posées.} Délais, taux d'absentéisme, pyramide de gravité aux
urgences sont aujourd'hui calibrés sur des sources publiées puis injectés. Sur des données réelles,
les contrôles qui vérifient que la chaîne retrouve le paramètre injecté n'auraient plus de second
membre, et devraient céder la place à des contrôles de forme, plus faibles et à écrire.

\textbf{Les relations ne sont plus injectées, et c'est le changement décisif : un résultat cesse de
démontrer et se met à établir.} Aujourd'hui, retrouver que l'absentéisme croît avec le délai à
l'intérieur d'une spécialité ne prouve rien d'autre que la capacité de la chaîne à retrouver ce
qu'on y a mis. Sur des données réelles, le même calcul deviendrait un constat sur le service — avec
tout ce que cela impose de prudence supplémentaire sur la population, le dénominateur et la période.

Trois conditions matérielles s'y ajoutent. Une convention d'accès aux données, que la contrainte de
confidentialité rend nécessaire et qui n'est pas du ressort d'un stage. Une extraction périodique
depuis le système d'origine, dont le format et la cadence restent à établir. Et les champs manquants nommés plus haut, qui
cesseraient d'être une limite documentée pour devenir un obstacle à lever.

\textbf{Apport personnel.} Le paragraphe qui suit relève de la rédaction personnelle et ne peut
être écrit que par l'auteur du travail.

\aRediger{apport-personnel}{ce que ce travail a apporté personnellement : les compétences acquises,
les difficultés rencontrées et la façon dont elles ont été surmontées, et ce que l'expérience du
service a changé dans la manière d'aborder un problème de
données}\releve{apport-personnel}


% La bibliographie liste ce que le document EMPLOIE. Un `\nocite{*}` imprimait auparavant
% toutes les entrées du fichier bibliographique, citées ou non, pour que la chaîne
% bibliographique soit vérifiable dès le squelette : un fichier cassé se serait vu à la
% première compilation plutôt qu'à la première citation. Ce motif a servi tant que le rapport
% n'était pas rédigé ; il ne sert plus, puisque vingt clés sont désormais citées et que toute
% rupture de la chaîne se voit aussitôt.
%
% Les entrées du registre des sources qu'aucun chapitre ne cite restent au registre avec leur
% motif de non-emploi. Elles n'ont pas disparu ; elles ne s'impriment plus.
\printbibliography[heading=bibintoc,title={Bibliographie}]

\appendix

\chapter{Dictionnaire de données}

Le tableau qui suit est produit mécaniquement depuis le registre des champs, et
n'est pas rédigé : un contrôle compare le fichier inclus ici au registre dont il
dérive.

% Fichier produit mécaniquement depuis le registre des champs : ne pas
% modifier à la main.

\subsection*{source.patients}
\begin{longtable}{p{2.5cm}p{1.8cm}p{2.8cm}p{1.3cm}p{2.2cm}p{5cm}}
\hline
Colonne & Type métier & Libellé Hosix & Provenance & Preuve & Note \\
\hline
\endhead
n\_ipp & code & N° IPP & OBS & REL-PAT.D01 & Relevé de champs, fiche patient, bloc Données principales. \\
nom & texte & Nom & OBS & REL-PAT.D02 & Relevé de champs, fiche patient, bloc Données principales. \\
nom\_famille\_1 & texte & Nom de famille 1 & OBS & REL-PAT.D03 & Relevé de champs, fiche patient, bloc Données principales. \\
nom\_famille\_2 & texte & Nom de famille 2 & OBS & REL-PAT.D04 & Relevé de champs, fiche patient, bloc Données principales. \\
sexe & code & Sexe & OBS & REL-PAT.D05 & Relevé de champs, fiche patient, bloc Données principales. \\
date\_naissance & date & D. Nai & OBS & REL-PAT.D06 & Relevé de champs, fiche patient, bloc Données principales. \\
type\_piece\_identite & code & Type pièce d'identité & OBS & REL-PAT.D09 & Relevé de champs, fiche patient, bloc Données principales. \\
n\_piece\_identite & texte & N° pièce d'identité & OBS & REL-PAT.D10 & Relevé de champs, fiche patient, bloc Données principales. \\
etat\_civil & code & E. Civil & OBS & REL-PAT.D11 & Relevé de champs, fiche patient, bloc Données principales. \\
type\_patient & code & Type patient & OBS & REL-PAT.D12 & Relevé de champs, fiche patient, bloc Données principales. \\
date\_photo & date & Date photo & OBS & REL-PAT.D13 & Relevé de champs, fiche patient, bloc Données principales. \\
modifie\_par & texte & Modifié par & OBS & REL-PAT.D14 & Relevé de champs, fiche patient, bloc Données principales. \\
cree\_par & texte & Créé par & OBS & REL-PAT.D15 & Relevé de champs, fiche patient, bloc Données principales. \\
date\_attribution & date & Date d'attribution & OBS & REL-PAT.D16 & Relevé de champs, fiche patient, bloc Données principales. \\
compagnie\_assurance & code & Compagnie d'assur. & OBS & REL-PAT.A01 & Relevé de champs, fiche patient, bloc Compagnie d'assurance. \\
police & texte & Police & OBS & REL-PAT.A02 & Relevé de champs, fiche patient, bloc Compagnie d'assurance. \\
n\_assure & texte & N° Assu & OBS & REL-PAT.A03 & Relevé de champs, fiche patient, bloc Compagnie d'assurance. \\
profession & texte & Profession & OBS & REL-PAT.A04 & Relevé de champs, fiche patient, bloc Compagnie d'assurance. \\
num\_inscription & texte & Num. inscription & OBS & REL-PAT.A05 & Relevé de champs, fiche patient, bloc Compagnie d'assurance. \\
date\_inscription & date & Date inscription & OBS & REL-PAT.A06 & Relevé de champs, fiche patient, bloc Compagnie d'assurance. \\
type\_domicile & code & Type & OBS & REL-PAT.H01 & Relevé de champs, fiche patient, bloc Domicile. \\
adresse & texte & Adresse & OBS & REL-PAT.H02 & Relevé de champs, fiche patient, bloc Domicile. \\
code\_postal & texte & Code postal & OBS & REL-PAT.H03 & Relevé de champs, fiche patient, bloc Domicile. \\
etat & code & État & OBS & REL-PAT.H04 & Relevé de champs, fiche patient, bloc Domicile. \\
ville & code & Ville & OBS & REL-PAT.H05 & Relevé de champs, fiche patient, bloc Domicile. \\
quartier & texte & Quartier & OBS & REL-PAT.H06 & Relevé de champs, fiche patient, bloc Domicile. \\
nationalite & code & Nationalité & OBS & REL-PAT.H07 & Relevé de champs, fiche patient, bloc Domicile. \\
telephone\_1 & texte & Téléphone 1 & OBS & REL-PAT.H08 & Relevé de champs, fiche patient, bloc Domicile. \\
telephone\_2 & texte & Téléphone 2 & OBS & REL-PAT.H09 & Relevé de champs, fiche patient, bloc Domicile. \\
telephone\_3 & texte & Téléphone 3 & OBS & REL-PAT.H10 & Relevé de champs, fiche patient, bloc Domicile. \\
telephone\_4 & texte & Téléphone 4 & OBS & REL-PAT.H11 & Relevé de champs, fiche patient, bloc Domicile. \\
avertissements\_sms & booleen & Avertissements SMS & OBS & REL-PAT.H12 & Relevé de champs, fiche patient, bloc Domicile. \\
email & texte & E-mail & OBS & REL-PAT.H13 & Relevé de champs, fiche patient, bloc Domicile. \\
avertissements\_email & booleen & Avertissements e-mail & OBS & REL-PAT.H14 & Relevé de champs, fiche patient, bloc Domicile. \\
environnement & code & Environnement & OBS & REL-PAT.H15 & Relevé de champs, fiche patient, bloc Domicile. \\
nom\_pere & texte & Nom. Père & OBS & REL-PAT.N01 & Relevé de champs, fiche patient, bloc Né. \\
nom\_mere & texte & Nom. Mère & OBS & REL-PAT.N02 & Relevé de champs, fiche patient, bloc Né. \\
etat\_naissance & code & Lieu de naissance - État & OBS & REL-PAT.N03 & Relevé de champs, fiche patient, bloc Né. \\
ville\_naissance & code & Ville & OBS & REL-PAT.N04 & Relevé de champs, fiche patient, bloc Né. \\
pays\_naissance & code & Pays & OBS & REL-PAT.N05 & Relevé de champs, fiche patient, bloc Né. \\
quartier\_naissance & texte & Quartier & OBS & REL-PAT.N06 & Relevé de champs, fiche patient, bloc Né. \\
commentaire & texte & Commentaire & OBS & REL-PAT.K01 & Relevé de champs, fiche patient, bloc Commentaire. \\
province & code & Province & OBS & REL-IPP.R09 & Relevé de champs, recherche d'IPP, colonne de résultat. La colonne n'apparaît pas au bloc Domicile de la fiche patient, où figurent État, Ville et Quartier ; elle est établie par la liste de résultats de la recherche, qui la restitue. \\
exitus & booleen & Exitus & OBS & REL-IPP.F04 & Relevé de champs, recherche d'IPP, filtre de population. L'existence d'un filtre sur les patients décédés établit qu'un indicateur de décès est stocké sur la fiche, bien qu'aucun champ de ce nom n'ait été relevé sur l'écran de la fiche patient. \\
date\_modification & horodatage & non\_releve & HYP & sans\_preuve\_externe & L'écran de la fiche patient porte un champ Modifié par sans horodatage adjacent, là où l'écran de rendez-vous appareille chaque agent à sa date. La colonne est posée parce que l'historisation de la dimension patient en type 2 exige un instant de changement. Si l'extraction réelle ne la fournissait pas, l'historisation se rabattrait sur la date d'extraction de la partition, à la granularité du jour au lieu de la seconde, et les changements survenus le même jour seraient confondus. \\
date\_extraction & date & non\_releve & HYP & sans\_preuve\_externe & Colonne technique de la zone d'atterrissage, absente du système observé. Elle porte la date de l'extraction dont la ligne provient et rend le rechargement d'une partition idempotent. Si une extraction réelle ne la fournissait pas, elle serait dérivée du nom du fichier partitionné, sans perte d'information. \\
\hline
\end{longtable}

\subsection*{source.rendez\_vous}
\begin{longtable}{p{2.5cm}p{1.8cm}p{2.8cm}p{1.3cm}p{2.2cm}p{5cm}}
\hline
Colonne & Type métier & Libellé Hosix & Provenance & Preuve & Note \\
\hline
\endhead
n\_rdv & code & non\_releve & HYP & sans\_preuve\_externe & Aucun identifiant de rendez-vous n'apparaît à l'écran Donner rendez-vous, dont la barre d'actions comporte pourtant Enregistrer, Supprimer et Chercher, trois opérations qui supposent une clé. La colonne est posée parce que la table des passages porte une référence de rendez-vous nullable, qui n'aurait pas de cible sans elle. Si le système en place clé le rendez-vous par un couple agenda et horodatage plutôt que par un numéro, la structure logique du modèle est inchangée et seule la forme de la clé diffère. \\
n\_ipp & code & N° IPP & OBS & REL-RDV.I01 & Relevé de champs, écran Donner rendez-vous, bloc Identification du patient. \\
agenda & code & Agenda & OBS & REL-RDV.R01 & Relevé de champs, écran Donner rendez-vous, bloc Rendez-vous. \\
activite & code & Activité & OBS & REL-RDV.R02 & Relevé de champs, écran Donner rendez-vous, bloc Rendez-vous. \\
origine & code & Origine & OBS & REL-RDV.R03 & Relevé de champs, écran Donner rendez-vous, bloc Rendez-vous. \\
hopital\_cs & code & Hôpital/C.S. & OBS & REL-RDV.R04 & Relevé de champs, écran Donner rendez-vous, bloc Rendez-vous. \\
medecin\_ext & texte & Médecin ext. & OBS & REL-RDV.R05 & Relevé de champs, écran Donner rendez-vous, bloc Rendez-vous. \\
service\_ext & texte & Service ext. & OBS & REL-RDV.R06 & Relevé de champs, écran Donner rendez-vous, bloc Rendez-vous. \\
observations & texte & Observations & OBS & REL-RDV.R07 & Relevé de champs, écran Donner rendez-vous, bloc Rendez-vous. \\
date\_rendez\_vous & horodatage & Date rendez-vous & OBS & REL-RDV.R08 & Relevé de champs, écran Donner rendez-vous, bloc Rendez-vous. \\
rdv\_supplementaire & booleen & Rendez-vous supplémentaire & OBS & REL-RDV.R09 & Relevé de champs, écran Donner rendez-vous, bloc Rendez-vous. \\
type\_attention & code & Type d'attention & OBS & REL-RDV.R10 & Relevé de champs, écran Donner rendez-vous, bloc Rendez-vous. \\
etat & code & État & OBS & REL-RDV.R11 & Relevé de champs, écran Donner rendez-vous, bloc Rendez-vous. \\
duree & entier & Durée & OBS & REL-RDV.R12 & Relevé de champs, écran Donner rendez-vous, bloc Rendez-vous. Durée exprimée en minutes, valeur nulle sur le rendez-vous observé. \\
date\_reception & horodatage & Date réception & OBS & REL-RDV.R13 & Relevé de champs, écran Donner rendez-vous, bloc Rendez-vous. \\
imprimer\_donnees & booleen & Imprimer données & OBS & REL-RDV.R14 & Relevé de champs, écran Donner rendez-vous, bloc Rendez-vous. \\
cree\_par & texte & Créé par & OBS & REL-RDV.C01 & Relevé de champs, écran Donner rendez-vous, bloc Contrôle de modifications. \\
date\_creation & horodatage & Date création & OBS & REL-RDV.C02 & Relevé de champs, écran Donner rendez-vous, bloc Contrôle de modifications. L'écart entre cette colonne et la date du rendez-vous est le délai d'obtention, grandeur centrale de l'analyse. Sur le rendez-vous observé, il valait une seconde. \\
modifie\_par & texte & Modifié par & OBS & REL-RDV.C03 & Relevé de champs, écran Donner rendez-vous, bloc Contrôle de modifications. \\
date\_mod & horodatage & Date mod. & OBS & REL-RDV.C04 & Relevé de champs, écran Donner rendez-vous, bloc Contrôle de modifications. \\
confirme\_par & texte & Confirmé par & OBS & REL-RDV.C05 & Relevé de champs, écran Donner rendez-vous, bloc Contrôle de modifications. Avec la colonne d'annulation, ce couple sépare une annulation déclarée d'une absence non prévenue ; sans lui, les deux se confondent. \\
date\_conf & horodatage & Date conf. & OBS & REL-RDV.C06 & Relevé de champs, écran Donner rendez-vous, bloc Contrôle de modifications. \\
annule\_par & texte & Annulé par & OBS & REL-RDV.C07 & Relevé de champs, écran Donner rendez-vous, bloc Contrôle de modifications. \\
date\_annul & horodatage & Date annul. & OBS & REL-RDV.C08 & Relevé de champs, écran Donner rendez-vous, bloc Contrôle de modifications. \\
liste\_attente\_service & code & Service & OBS & REL-RDV.L01 & Relevé de champs, écran Donner rendez-vous, bloc Liste d'attente des consultations. \\
liste\_attente\_agenda & code & Agenda & OBS & REL-RDV.L02 & Relevé de champs, écran Donner rendez-vous, bloc Liste d'attente des consultations. \\
liste\_attente\_activite & code & Activité & OBS & REL-RDV.L03 & Relevé de champs, écran Donner rendez-vous, bloc Liste d'attente des consultations. \\
date\_extraction & date & non\_releve & HYP & sans\_preuve\_externe & Colonne technique de la zone d'atterrissage, absente du système observé. Elle porte la date de l'extraction dont la ligne provient et rend le rechargement d'une partition idempotent. Si une extraction réelle ne la fournissait pas, elle serait dérivée du nom du fichier partitionné, sans perte d'information. \\
\hline
\end{longtable}

\subsection*{source.passages}
\begin{longtable}{p{2.5cm}p{1.8cm}p{2.8cm}p{1.3cm}p{2.2cm}p{5cm}}
\hline
Colonne & Type métier & Libellé Hosix & Provenance & Preuve & Note \\
\hline
\endhead
n\_passage & code & non\_releve & DOC & S-08 & Le module Consultas de l'éditeur documente la gestion des consultations et la production de leurs listings, ce qui suppose un identifiant de passage. \\
n\_ipp & code & N° IPP & OBS & REL-PAT.D01 & Relevé de champs, fiche patient, bloc Données principales. \\
type\_passage & code & non\_releve & DOC & S-06 & La liste de travail du module Médicos regroupe les patients par type d'épisode : hospitalisés, consultation, urgences. Le système type donc ses épisodes. \\
service & code & non\_releve & DOC & S-27 & Article 27 du règlement intérieur : l'hôpital s'organise en services, et l'activité s'y rattache. \\
activite & code & Activité & OBS & REL-RDV.R02 & Relevé de champs, écran Donner rendez-vous, bloc Rendez-vous. Le même référentiel d'activités qualifie le rendez-vous et le passage qui l'honore. \\
n\_rdv & code & non\_releve & HYP & sans\_preuve\_externe & La colonne rattache le passage au rendez-vous qui l'a programmé, et hérite du caractère hypothétique de l'identifiant de rendez-vous lui-même, qui n'apparaît sur aucun écran. Elle est nullable : un passage peut survenir sans rendez-vous préalable. Si le système en place ne conservait pas ce rattachement, le délai d'obtention resterait mesurable sur la table des rendez-vous, mais le taux d'honoration des rendez-vous ne le serait plus. \\
mode\_prise\_en\_charge & code & non\_releve & DOC & S-18 & Le tiers payant distingue les modes de prise en charge selon le régime et le caractère ambulatoire ou hospitalier de la prestation. \\
date\_heure\_entree & horodatage & non\_releve & DOC & S-08 & Le module Consultas gère la programmation et le déroulement de la consultation, dont l'heure d'entrée. \\
date\_heure\_sortie & horodatage & non\_releve & DOC & S-08 & Idem, borne de fin du passage. \\
medecin & texte & non\_releve & DOC & S-06 & Le module Médicos rattache l'épisode au praticien qui le prend en charge. \\
cree\_par & texte & Créé par & OBS & REL-RDV.C01 & Relevé de champs, écran Donner rendez-vous, bloc Contrôle de modifications. Le même bloc d'audit est reproduit d'un écran à l'autre. \\
date\_creation & horodatage & Date création & OBS & REL-RDV.C02 & Relevé de champs, écran Donner rendez-vous, bloc Contrôle de modifications. Le même bloc d'audit est reproduit d'un écran à l'autre. \\
date\_extraction & date & non\_releve & HYP & sans\_preuve\_externe & Colonne technique de la zone d'atterrissage, absente du système observé. Elle porte la date de l'extraction dont la ligne provient et rend le rechargement d'une partition idempotent. Si une extraction réelle ne la fournissait pas, elle serait dérivée du nom du fichier partitionné, sans perte d'information. \\
\hline
\end{longtable}

\subsection*{source.mouvements}
\begin{longtable}{p{2.5cm}p{1.8cm}p{2.8cm}p{1.3cm}p{2.2cm}p{5cm}}
\hline
Colonne & Type métier & Libellé Hosix & Provenance & Preuve & Note \\
\hline
\endhead
n\_sejour & code & non\_releve & DOC & S-27 & Article 40 : les formalités d'admission ouvrent un dossier de séjour identifié. \\
n\_ipp & code & N° IPP & OBS & REL-PAT.D01 & Relevé de champs, fiche patient, bloc Données principales. \\
date\_heure\_admission & horodatage & non\_releve & DOC & S-27 & Article 40, formalités d'admission ordinaire. \\
mode\_admission & code & non\_releve & DOC & S-27 & Articles 40, 42 et 45 : l'admission est ordinaire ou en urgence, et les formalités diffèrent. \\
service\_accueil & code & non\_releve & DOC & S-27 & Article 44 : l'hospitalisation est ordonnée vers un service. \\
lit & texte & non\_releve & HYP & sans\_preuve\_externe & Aucune source ne documente l'identification du lit occupé, alors que la capacité litière fonctionnelle est publiée par établissement et que l'article 42 impose l'admission même en cas d'indisponibilité de lits, ce qui suppose un suivi de l'occupation. La colonne est posée parce que le taux d'occupation du tableau de bord se recalcule depuis les journées d'hospitalisation, non depuis les lits ; si elle était fausse ou absente, aucun indicateur du chapitre sur l'activité ne changerait. \\
n\_mutation & code & non\_releve & DOC & S-27 & Article 35 : le service gère les effectifs des patients et leurs mouvements à l'intérieur de l'hôpital. \\
service\_origine & code & non\_releve & DOC & S-27 & Article 35, mouvements internes. \\
service\_destination & code & non\_releve & DOC & S-27 & Article 35, mouvements internes. \\
date\_heure\_mutation & horodatage & non\_releve & DOC & S-27 & Article 35, mouvements internes. \\
date\_heure\_sortie & horodatage & non\_releve & DOC & S-27 & Article 79 : les formalités de sortie closent le séjour. \\
mode\_sortie & code & non\_releve & DOC & S-27 & Articles 79 et 80 : sortie régulière avec billet, sortie à l'insu du personnel, décès. \\
date\_extraction & date & non\_releve & HYP & sans\_preuve\_externe & Colonne technique de la zone d'atterrissage, absente du système observé. Elle porte la date de l'extraction dont la ligne provient et rend le rechargement d'une partition idempotent. Si une extraction réelle ne la fournissait pas, elle serait dérivée du nom du fichier partitionné, sans perte d'information. \\
\hline
\end{longtable}

\subsection*{source.prises\_en\_charge}
\begin{longtable}{p{2.5cm}p{1.8cm}p{2.8cm}p{1.3cm}p{2.2cm}p{5cm}}
\hline
Colonne & Type métier & Libellé Hosix & Provenance & Preuve & Note \\
\hline
\endhead
n\_prise\_en\_charge & code & non\_releve & DOC & S-27 & Article 79 : les formalités de sortie comportent la signature des documents de prise en charge, document identifié. \\
n\_ipp & code & N° IPP & OBS & REL-PAT.D01 & Relevé de champs, fiche patient, bloc Données principales. \\
n\_episode & code & non\_releve & DOC & S-27 & Article 79 : la prise en charge se rattache à l'épisode qu'elle couvre. \\
type\_episode & code & non\_releve & DOC & S-27 & Article 36 : taxonomie réglementaire des modes d'utilisation de l'hôpital. \\
organisme & code & non\_releve & DOC & S-15 & Les régimes de l'assurance maladie obligatoire sont distincts et l'organisme gestionnaire diffère de l'un à l'autre. \\
n\_assure & texte & N° Assu & OBS & REL-PAT.A03 & Relevé de champs, fiche patient, bloc Compagnie d'assurance. \\
date\_verification & horodatage & non\_releve & DOC & S-27 & Article 40 : le patient présente les documents exigés selon son statut de couverture, ce qui date la vérification. \\
etat & code & non\_releve & DOC & S-19 & Aucune demande de prise en charge préalable n'est exigée dans le secteur public : l'état constate l'ouverture des droits, il n'enregistre pas un accord. Valeurs retenues : droits ouverts, droits fermés, non vérifié. \\
taux\_prise\_en\_charge & decimal & non\_releve & DOC & S-18 & Cent pour cent de la tarification nationale de référence en hospitalisation et en hôpital de jour, quatre-vingts pour cent en ambulatoire au-delà de deux cents dirhams. \\
date\_extraction & date & non\_releve & HYP & sans\_preuve\_externe & Colonne technique de la zone d'atterrissage, absente du système observé. Elle porte la date de l'extraction dont la ligne provient et rend le rechargement d'une partition idempotent. Si une extraction réelle ne la fournissait pas, elle serait dérivée du nom du fichier partitionné, sans perte d'information. \\
\hline
\end{longtable}

\subsection*{source.factures}
\begin{longtable}{p{2.5cm}p{1.8cm}p{2.8cm}p{1.3cm}p{2.2cm}p{5cm}}
\hline
Colonne & Type métier & Libellé Hosix & Provenance & Preuve & Note \\
\hline
\endhead
n\_facture & code & non\_releve & DOC & S-09 & Le module Facturación de Hosix Core documente l'émission de factures et de listings personnalisés, ce qui suppose un identifiant de facture. \\
n\_ipp & code & N° IPP & OBS & REL-PAT.D01 & Relevé de champs, fiche patient, bloc Données principales. \\
n\_episode & code & non\_releve & DOC & S-09 & Le module Facturación récupère dans chacun des autres modules l'information nécessaire à la facturation en captant les mouvements générés, ce qui suppose un identifiant d'épisode rattaché. \\
type\_episode & code & non\_releve & DOC & S-27 & L'article 36 du règlement intérieur des hôpitaux établit une taxonomie réglementaire des modes d'utilisation de l'hôpital, que la facture doit porter. \\
code\_diagnostic\_cim10 & code & non\_releve & DOC & S-27 & L'article 35 du règlement intérieur impose d'établir la facturation sur la base de la classification des maladies, soit la CIM-10. \\
date\_facture & date & non\_releve & DOC & S-09 & Le module Facturación de l'éditeur documente l'émission de factures, ce qui suppose une date d'émission. \\
type\_facture & code & non\_releve & DOC & S-09 & Le module Facturación admet la facturation aux entités publiques, aux entités privées et particuliers, et aux mutuelles d'accidents. \\
service\_emetteur & code & non\_releve & DOC & S-27 & L'article 35 du règlement intérieur charge le service d'accueil et d'admission d'établir la facturation, ce qui rattache chaque facture à un service émetteur. \\
etat & code & non\_releve & DOC & S-09 & Le module Facturación contrôle explicitement l'état des activités réalisées : facturées, en attente de facturation, et autres états. \\
montant\_total & decimal & non\_releve & DOC & S-17 & Le montant se calcule sur la grille des lettres clés de la tarification nationale de référence du secteur public. \\
part\_organisme & decimal & non\_releve & DOC & S-18 & Les taux de prise en charge de la tarification nationale de référence fixent la part organisme à 100 \% en hospitalisation et hôpital de jour, et à 80 \% en ambulatoire au-delà de 200 dirhams. \\
part\_patient & decimal & non\_releve & DOC & S-18 & La part patient est le complément de la part organisme, sauf pour les bénéficiaires du régime AMO-Tadamon en structure publique où l'État prend en charge le ticket modérateur et où la part patient devient nulle, comme l'établit une seconde source (S-15). \\
cree\_par & texte & Créé par & OBS & REL-RDV.C01 & Relevé de champs, écran Donner rendez-vous, bloc Contrôle de modifications. Le même bloc d'audit est reproduit d'un écran à l'autre. \\
date\_creation & horodatage & Date création & OBS & REL-RDV.C02 & Relevé de champs, écran Donner rendez-vous, bloc Contrôle de modifications. Le même bloc d'audit est reproduit d'un écran à l'autre. \\
date\_extraction & date & non\_releve & HYP & sans\_preuve\_externe & Colonne technique de la zone d'atterrissage, absente du système observé. Elle porte la date de l'extraction dont la ligne provient et rend le rechargement d'une partition idempotent. Si une extraction réelle ne la fournissait pas, elle serait dérivée du nom du fichier partitionné, sans perte d'information. \\
\hline
\end{longtable}

\subsection*{source.lignes\_facture}
\begin{longtable}{p{2.5cm}p{1.8cm}p{2.8cm}p{1.3cm}p{2.2cm}p{5cm}}
\hline
Colonne & Type métier & Libellé Hosix & Provenance & Preuve & Note \\
\hline
\endhead
n\_facture & code & non\_releve & DOC & S-09 & Le module Facturación de l'éditeur documente l'émission de factures composées de lignes. \\
n\_ligne & entier & non\_releve & DOC & S-09 & Le module Facturación produit une documentation et des listings personnalisés composés de lignes détaillées. \\
code\_acte & code & non\_releve & DOC & S-27 & L'article 35 du règlement intérieur impose une facturation sur la base des nomenclatures des actes, et la grille des lettres clés du secteur public rattache un code à chaque acte. \\
libelle\_acte & texte & non\_releve & DOC & S-17 & La grille des lettres clés de la tarification nationale de référence du secteur public porte un libellé pour chaque acte, en complément de la nomenclature imposée par l'article 35 du règlement intérieur. \\
lettre\_cle & code & non\_releve & DOC & S-17 & La tarification nationale de référence rémunère l'acte par une lettre clé et un coefficient. \\
coefficient & decimal & non\_releve & DOC & S-17 & La tarification nationale de référence rémunère l'acte par une lettre clé et un coefficient. \\
quantite & entier & non\_releve & HYP & sans\_preuve\_externe & La structure d'une ligne de facturation hospitalière comporte usuellement une quantité distincte du coefficient de la lettre clé, mais aucune des sources ouvertes ne la documente pour le secteur public marocain. La colonne est posée parce que sans elle un acte répété le même jour se confondrait avec un acte unique. Si elle était absente du système réel, les montants resteraient exacts et seul le dénombrement des actes par ligne serait affecté. \\
tarif\_unitaire & decimal & non\_releve & DOC & S-17 & Le tarif unitaire se lit sur la grille des lettres clés de la tarification nationale de référence. \\
montant & decimal & non\_releve & DOC & S-17 & Le montant de la ligne se calcule sur la grille des lettres clés de la tarification nationale de référence. \\
service\_executant & code & non\_releve & DOC & S-09 & Le module Facturación rattache chaque activité facturable au service qui l'a exécutée. \\
date\_acte & date & non\_releve & DOC & S-27 & L'article 42 du règlement intérieur dispose qu'aux urgences la facturation n'est entamée qu'après l'engagement de la prise en charge médicale, ce qui date l'acte facturé. \\
date\_extraction & date & non\_releve & HYP & sans\_preuve\_externe & Colonne technique de la zone d'atterrissage, absente du système observé. Elle porte la date de l'extraction dont la ligne provient et rend le rechargement d'une partition idempotent. Si une extraction réelle ne la fournissait pas, elle serait dérivée du nom du fichier partitionné, sans perte d'information. \\
\hline
\end{longtable}

\subsection*{source.passages\_urgences}
\begin{longtable}{p{2.5cm}p{1.8cm}p{2.8cm}p{1.3cm}p{2.2cm}p{5cm}}
\hline
Colonne & Type métier & Libellé Hosix & Provenance & Preuve & Note \\
\hline
\endhead
n\_passage & code & non\_releve & DOC & S-27 & L'article 45 du règlement intérieur prescrit que les formalités d'admission en urgence sont enregistrées, ce qui identifie chaque passage. \\
n\_ipp & code & N° IPP & OBS & REL-PAT.D01 & Relevé de champs, fiche patient, bloc Données principales. \\
date\_heure\_arrivee & horodatage & non\_releve & DOC & S-27 & L'article 42 du règlement intérieur fait de l'accueil aux urgences l'acte fondateur du passage. \\
mode\_arrivee & code & non\_releve & DOC & S-12 & Le mode d'arrivée aux urgences est documenté par la protection civile, l'ambulance SMUR, les moyens propres et le transport privé. \\
motif\_recours & code & non\_releve & DOC & S-13 & L'étude sur le recours non approprié aux urgences classe les motifs de consultation par chapitre de la CIM-10. \\
niveau\_tri & code & non\_releve & DOC & S-12 & Les chiffres du ministère répartissent les passages entre urgences vitales, environ 10 \%, et consultations médicales non urgentes, 64 \%, le solde de 26 \% correspondant aux urgences réelles non vitales. Réserve : cette source ne distingue que trois groupes de gravité, tandis que le modèle retient cinq niveaux ; la colonne existe, c'est son échelle qui reste ouverte. \\
date\_heure\_pec\_medicale & horodatage & non\_releve & DOC & S-27 & L'article 42 du règlement intérieur dispose que la procédure de facturation n'est entamée qu'après l'engagement de la prise en charge médicale, ce qui date cet engagement. \\
date\_heure\_sortie & horodatage & non\_releve & DOC & S-27 & L'article 43 du règlement intérieur impose au patient admis aux urgences de s'acquitter des frais avant sa sortie, ce que le règlement précède. \\
orientation\_sortie & code & non\_releve & DOC & S-27 & Les articles 42, 44, 46 et 47 du règlement intérieur distinguent le retour à domicile, l'hospitalisation, le transfert, la sortie contre avis médical et le décès comme issues du passage. \\
service\_orientation & code & non\_releve & DOC & S-27 & L'article 44 du règlement intérieur dispose que l'hospitalisation d'urgence est ordonnée vers un service. \\
motif\_transfert & code & non\_releve & DOC & S-27 & L'article 47 du règlement intérieur prévoit le transfert lorsque les soins requis relèvent d'une discipline ou d'une technique n'existant pas à l'hôpital. \\
consentement\_transfert & booleen & non\_releve & DOC & S-27 & L'article 47 du règlement intérieur exige le consentement écrit du patient pour un transfert, sauf extrême urgence. \\
famille\_informee & booleen & non\_releve & DOC & S-27 & Les articles 47 et 67 du règlement intérieur imposent d'informer la famille, lors d'un transfert comme pour un patient mineur, incapable ou inconscient. \\
inventaire\_effets & booleen & non\_releve & DOC & S-27 & L'article 45 du règlement intérieur impose un inventaire contradictoire des effets personnels lorsque le patient est inconscient. \\
date\_extraction & date & non\_releve & HYP & sans\_preuve\_externe & Colonne technique de la zone d'atterrissage, absente du système observé. Elle porte la date de l'extraction dont la ligne provient et rend le rechargement d'une partition idempotent. Si une extraction réelle ne la fournissait pas, elle serait dérivée du nom du fichier partitionné, sans perte d'information. \\
\hline
\end{longtable}

\subsection*{source.encaissements}
\begin{longtable}{p{2.5cm}p{1.8cm}p{2.8cm}p{1.3cm}p{2.2cm}p{5cm}}
\hline
Colonne & Type métier & Libellé Hosix & Provenance & Preuve & Note \\
\hline
\endhead
n\_encaissement & code & non\_releve & DOC & S-27 & L'article 79 du règlement intérieur range le règlement des frais parmi les formalités de sortie, ce qui identifie chaque encaissement. \\
n\_facture & code & non\_releve & DOC & S-27 & L'article 79 du règlement intérieur rattache le règlement des frais à la facture émise. \\
date\_encaissement & horodatage & non\_releve & DOC & S-27 & L'article 79 du règlement intérieur date le règlement des frais parmi les formalités de sortie. \\
mode\_reglement & code & non\_releve & HYP & sans\_preuve\_externe & Aucune source ouverte ne documente les moyens de paiement acceptés au guichet d'un hôpital public marocain ; la Cour des comptes décrit seulement un versement direct chez le régisseur. La colonne est posée parce que la page de recouvrement du tableau de bord distingue les canaux d'encaissement. Si la répartition retenue était fausse, aucun taux de recouvrement ni aucune ancienneté de créance n'en serait affecté, ces grandeurs ne dépendant que du montant et de la date. \\
montant & decimal & non\_releve & DOC & S-18 & Le montant encaissé suit les taux de prise en charge de la tarification nationale de référence. \\
regisseur & texte & non\_releve & DOC & S-20 & Le rapport de la Cour des comptes décrit un paiement direct chez le régisseur. \\
billet\_sortie\_delivre & booleen & non\_releve & DOC & S-27 & L'article 79 du règlement intérieur conditionne la délivrance du billet de sortie au règlement des frais ou à la signature des documents de prise en charge. \\
date\_extraction & date & non\_releve & HYP & sans\_preuve\_externe & Colonne technique de la zone d'atterrissage, absente du système observé. Elle porte la date de l'extraction dont la ligne provient et rend le rechargement d'une partition idempotent. Si une extraction réelle ne la fournissait pas, elle serait dérivée du nom du fichier partitionné, sans perte d'information. \\
\hline
\end{longtable}

\subsection*{source.creances}
\begin{longtable}{p{2.5cm}p{1.8cm}p{2.8cm}p{1.3cm}p{2.2cm}p{5cm}}
\hline
Colonne & Type métier & Libellé Hosix & Provenance & Preuve & Note \\
\hline
\endhead
n\_creance & code & non\_releve & DOC & S-27 & L'article 9, paragraphe b, du règlement intérieur charge le pôle des affaires administratives du recouvrement des créances de l'établissement, ce qui identifie chaque créance. \\
n\_facture & code & non\_releve & DOC & S-27 & L'article 9, paragraphe b, du règlement intérieur rattache la créance recouvrée à la facture qui l'a fait naître. \\
date\_naissance\_creance & date & non\_releve & DOC & S-27 & L'article 79 du règlement intérieur fait naître la créance aux formalités de sortie, lorsque la facturation n'est pas soldée. \\
montant\_du & decimal & non\_releve & DOC & S-20 & Le rapport de la Cour des comptes établit l'existence d'un état des prestations non recouvrées, dont le montant dû. \\
montant\_recouvre & decimal & non\_releve & DOC & S-20 & Le rapport de la Cour des comptes établit l'existence d'un état des prestations non recouvrées, dont le montant recouvré. \\
montant\_restant & decimal & non\_releve & DOC & S-20 & Le rapport de la Cour des comptes établit l'existence d'un état des prestations non recouvrées, dont le montant restant dû. Cette colonne est conservée bien qu'elle se déduise des deux colonnes voisines : sa redondance interne à la ligne ne dépend pas de la date de lecture, et l'écart éventuel entre les trois montants est un contrôle de cohérence exploitable. \\
type\_debiteur & code & non\_releve & DOC & S-09 & Le module Facturación distingue les débiteurs publics, privés, mutualistes et particuliers. \\
motif\_non\_recouvrement & code & non\_releve & DOC & S-20 & Le rapport de la Cour des comptes relève des sorties sans règlement, des patients non identifiés et des évasions comme motifs de non-recouvrement. \\
date\_extraction & date & non\_releve & HYP & sans\_preuve\_externe & Colonne technique de la zone d'atterrissage, absente du système observé. Elle porte la date de l'extraction dont la ligne provient et rend le rechargement d'une partition idempotent. Si une extraction réelle ne la fournissait pas, elle serait dérivée du nom du fichier partitionné, sans perte d'information. \\
\hline
\end{longtable}

\subsection*{source.relances}
\begin{longtable}{p{2.5cm}p{1.8cm}p{2.8cm}p{1.3cm}p{2.2cm}p{5cm}}
\hline
Colonne & Type métier & Libellé Hosix & Provenance & Preuve & Note \\
\hline
\endhead
n\_relance & code & non\_releve & HYP & sans\_preuve\_externe & Aucune source ouverte ne documente une procédure de relance formalisée dans cet établissement ; le rapport de la Cour des comptes constate au contraire l'absence de diligences de recouvrement. La colonne identifie chaque relance parce que la page de recouvrement du tableau de bord la mobilise ; ce que la chaîne démontre ici est la calculabilité de l'indicateur, non une mesure. \\
n\_creance & code & non\_releve & HYP & sans\_preuve\_externe & Aucune source ouverte ne documente une procédure de relance formalisée ; la seule source disponible sur le recouvrement décrit au contraire son absence. La colonne rattache la relance à la créance parce que le tableau de bord en a besoin pour sa page recouvrement ; elle démontre une calculabilité, non une pratique observée. \\
date\_relance & date & non\_releve & HYP & sans\_preuve\_externe & Aucune source ouverte ne documente de procédure de relance formalisée dans cet établissement, et le rapport de la Cour des comptes constate précisément l'absence de diligences de recouvrement. La date de relance est posée parce que la page de recouvrement du tableau de bord en a besoin ; elle atteste une calculabilité, non un fait mesuré. \\
canal & code & non\_releve & HYP & sans\_preuve\_externe & Aucune source ouverte ne documente les canaux de relance employés dans cet établissement ; la source disponible sur le recouvrement décrit au contraire une carence de diligences. La colonne existe parce que le tableau de bord distingue les canaux à sa page recouvrement ; elle démontre une capacité de calcul, pas une observation. \\
resultat & code & non\_releve & HYP & sans\_preuve\_externe & Aucune source ouverte ne documente le résultat d'une procédure de relance dans cet établissement, dont la seule source disponible sur le recouvrement décrit l'absence de diligences. La colonne est posée parce que la page de recouvrement du tableau de bord en a besoin ; ce qu'elle démontre est une calculabilité, non une mesure. \\
date\_extraction & date & non\_releve & HYP & sans\_preuve\_externe & Colonne technique de la zone d'atterrissage, absente du système observé. Elle porte la date de l'extraction dont la ligne provient et rend le rechargement d'une partition idempotent. Si une extraction réelle ne la fournissait pas, elle serait dérivée du nom du fichier partitionné, sans perte d'information. \\
\hline
\end{longtable}


% Relevé exhaustif des quatre écrans observés.
%
% CE FICHIER EST UNE ANNEXE, ET IL EST LU PAR LE CONTRÔLE DE PROVENANCE au même titre qu'un
% fichier de chapitre : `tests/test_provenance_des_chapitres.py` porte la liste déclarative des
% fichiers d'annexe qu'il examine, et cette liste est confrontée aux inclusions de
% `report/annexes.tex` dans les deux sens.
%
% Les quatre tableaux qui suivent sont reproduits à l'identique de ce que portait le chapitre~2
% avant que le corps ne soit allégé. Aucun élément n'a été retiré, aucun libellé retouché : le
% corps en donne désormais la synthèse et un extrait, l'annexe la totalité.
%
% Ce fichier repose sur :
%   DOC: (aucun)
%   OBS: REL-DEM.I01, REL-DEM.I02, REL-DEM.P01, REL-DEM.P02, REL-DEM.P03, REL-DEM.P04, REL-DEM.P05,
%   OBS: REL-DEM.M01, REL-DEM.M02, REL-DEM.M03, REL-DEM.M04, REL-IPP.O01, REL-IPP.O02, REL-IPP.O03,
%   OBS: REL-IPP.O04, REL-IPP.C01, REL-IPP.C02, REL-IPP.C03, REL-IPP.C04, REL-IPP.C05, REL-IPP.C06,
%   OBS: REL-IPP.C07, REL-IPP.C08, REL-IPP.C09, REL-IPP.C10, REL-IPP.C11, REL-IPP.C12, REL-IPP.C13,
%   OBS: REL-IPP.F01, REL-IPP.F02, REL-IPP.F03, REL-IPP.F04, REL-IPP.R01, REL-IPP.R02, REL-IPP.R03,
%   OBS: REL-IPP.R04, REL-IPP.R05, REL-IPP.R06, REL-IPP.R07, REL-IPP.R08, REL-IPP.R09, REL-IPP.R10,
%   OBS: REL-IPP.R11, REL-IPP.R12, REL-PAT.T01, REL-PAT.T02, REL-PAT.T03, REL-PAT.D01, REL-PAT.D02,
%   OBS: REL-PAT.D03, REL-PAT.D04, REL-PAT.D05, REL-PAT.D06, REL-PAT.D07, REL-PAT.D08, REL-PAT.D09,
%   OBS: REL-PAT.D10, REL-PAT.D11, REL-PAT.D12, REL-PAT.D13, REL-PAT.D14, REL-PAT.D15, REL-PAT.D16,
%   OBS: REL-PAT.A01, REL-PAT.A02, REL-PAT.A03, REL-PAT.A04, REL-PAT.A05, REL-PAT.A06, REL-PAT.H01,
%   OBS: REL-PAT.H02, REL-PAT.H03, REL-PAT.H04, REL-PAT.H05, REL-PAT.H06, REL-PAT.H07, REL-PAT.H08,
%   OBS: REL-PAT.H09, REL-PAT.H10, REL-PAT.H11, REL-PAT.H12, REL-PAT.H13, REL-PAT.H14, REL-PAT.H15,
%   OBS: REL-PAT.N01, REL-PAT.N02, REL-PAT.N03, REL-PAT.N04, REL-PAT.N05, REL-PAT.N06, REL-PAT.K01,
%   OBS: REL-RDV.I01, REL-RDV.I02, REL-RDV.I03, REL-RDV.I04, REL-RDV.I05, REL-RDV.I06, REL-RDV.I07,
%   OBS: REL-RDV.I08, REL-RDV.I09, REL-RDV.I10, REL-RDV.R01, REL-RDV.R02, REL-RDV.R03, REL-RDV.R04,
%   OBS: REL-RDV.R05, REL-RDV.R06, REL-RDV.R07, REL-RDV.R08, REL-RDV.R09, REL-RDV.R10, REL-RDV.R11,
%   OBS: REL-RDV.R12, REL-RDV.R13, REL-RDV.R14, REL-RDV.C01, REL-RDV.C02, REL-RDV.C03, REL-RDV.C04,
%   OBS: REL-RDV.C05, REL-RDV.C06, REL-RDV.C07, REL-RDV.C08, REL-RDV.L01, REL-RDV.L02, REL-RDV.L03
%   HYP: (aucun)
%   CHI: (aucun)
%   SER: (aucune)
%   SECTIONS: 4

\chapitreDAnnexe{Relevé des écrans observés}\label{annexe:releve}

Ces quatre tableaux restituent l'intégralité du relevé conduit au poste, élément par élément. Le
chapitre~\ref{chap:systeme} en donne la synthèse et un extrait ; ce qui suit en donne la totalité.

Chaque élément porte un identifiant stable, le bloc de l'écran auquel il appartient, son libellé
reproduit à l'identique, son type apparent, et la valeur affichée lorsqu'une valeur l'était. Les
boutons de barre d'outils ne figurent pas à ces tableaux : ils déclenchent une opération et ne
portent aucune donnée, et le relevé les déclare comme tels.

\section{Écran de démarrage}

Onze éléments, en trois blocs.

\begin{longtable}{@{}p{2.9cm}p{3.9cm}p{2.8cm}p{4.0cm}@{}}
\toprule
Identifiant & Libellé à l'écran & Type apparent & Valeur observée \\
\midrule\endfirsthead
\toprule
Identifiant & Libellé à l'écran & Type apparent & Valeur observée \\
\midrule\endhead
\multicolumn{4}{@{}l}{\textbf{Identification du site}} \\
\texttt{REL-DEM.I01} & Identification du site & texte & \texttt{005 Fès-Meknès; Sidi Said} \\
\texttt{REL-DEM.I02} & Ministère de tutelle & texte & \texttt{Ministère de la Santé} \\
\addlinespace
\multicolumn{4}{@{}l}{\textbf{Profils applicatifs}} \\
\texttt{REL-DEM.P01} & MSM - FACTURATION SAA & liste & — \\
\texttt{REL-DEM.P02} & MSM - GESTION DE RDV & liste & — \\
\texttt{REL-DEM.P03} & MSM - RECOUVREMENT & liste & — \\
\texttt{REL-DEM.P04} & MSM - URGENCE & liste & — \\
\texttt{REL-DEM.P05} & MSM RAPPORTS ET STATISTIQUES & liste & — \\
\addlinespace
\multicolumn{4}{@{}l}{\textbf{Menus principaux}} \\
\texttt{REL-DEM.M01} & Hosix.NET & menu & — \\
\texttt{REL-DEM.M02} & Données Patient & menu & — \\
\texttt{REL-DEM.M03} & Centre de Consultation & menu & — \\
\texttt{REL-DEM.M04} & Factures & menu & — \\
\addlinespace
\bottomrule
\end{longtable}

\section{Recherche d'identifiant patient}

Trente-trois éléments, en quatre blocs.

\begin{longtable}{@{}p{2.9cm}p{3.9cm}p{2.8cm}p{4.0cm}@{}}
\toprule
Identifiant & Libellé à l'écran & Type apparent & Valeur observée \\
\midrule\endfirsthead
\toprule
Identifiant & Libellé à l'écran & Type apparent & Valeur observée \\
\midrule\endhead
\multicolumn{4}{@{}l}{\textbf{Onglets}} \\
\texttt{REL-IPP.O01} & Chercher & onglet & — \\
\texttt{REL-IPP.O02} & Chercher MPI & onglet & — \\
\texttt{REL-IPP.O03} & Chercher MPI Pondéré & onglet & — \\
\texttt{REL-IPP.O04} & Nouvel IPP & onglet & — \\
\addlinespace
\multicolumn{4}{@{}l}{\textbf{Critères de recherche}} \\
\texttt{REL-IPP.C01} & N° document & texte & \texttt{C.I.N} \\
\texttt{REL-IPP.C02} & Nom & texte & — \\
\texttt{REL-IPP.C03} & Prénom & texte & — \\
\texttt{REL-IPP.C04} & Date de naissance & date & — \\
\texttt{REL-IPP.C05} & Mois & entier & — \\
\texttt{REL-IPP.C06} & Année & entier & — \\
\texttt{REL-IPP.C07} & Approximation en années & entier & — \\
\texttt{REL-IPP.C08} & Âge & entier & — \\
\texttt{REL-IPP.C09} & Téléphone & texte & — \\
\texttt{REL-IPP.C10} & N° document de contact & texte & — \\
\texttt{REL-IPP.C11} & Sexe & liste & — \\
\texttt{REL-IPP.C12} & État civil & liste & — \\
\texttt{REL-IPP.C13} & N° assuré & texte & — \\
\addlinespace
\multicolumn{4}{@{}l}{\textbf{Filtres de population}} \\
\texttt{REL-IPP.F01} & Patients Hospitalisés & option & — \\
\texttt{REL-IPP.F02} & Patients aux Urgences & option & — \\
\texttt{REL-IPP.F03} & Tous les patients & option & — \\
\texttt{REL-IPP.F04} & Exitus & case\ a\ cocher & — \\
\addlinespace
\multicolumn{4}{@{}l}{\textbf{Colonnes de résultat}} \\
\texttt{REL-IPP.R01} & IPP & colonne & — \\
\texttt{REL-IPP.R02} & Nom & colonne & — \\
\texttt{REL-IPP.R03} & D. Naissance & colonne & — \\
\texttt{REL-IPP.R04} & N° pièce d'identité & colonne & — \\
\texttt{REL-IPP.R05} & N° Assu & colonne & — \\
\texttt{REL-IPP.R06} & E. Civil & colonne & — \\
\texttt{REL-IPP.R07} & Tél & colonne & — \\
\texttt{REL-IPP.R08} & Sexe & colonne & — \\
\texttt{REL-IPP.R09} & Province & colonne & — \\
\texttt{REL-IPP.R10} & Ville & colonne & — \\
\texttt{REL-IPP.R11} & Code postal & colonne & — \\
\texttt{REL-IPP.R12} & Probabilité & colonne & — \\
\addlinespace
\bottomrule
\end{longtable}

\section{Fiche patient}

Cinquante-sept éléments, dont dix boutons de barre d'outils hors tableau. Les quarante-sept
restants se répartissent en trois onglets et cinq blocs de données.

\begin{longtable}{@{}p{2.9cm}p{3.9cm}p{2.8cm}p{4.0cm}@{}}
\toprule
Identifiant & Libellé à l'écran & Type apparent & Valeur observée \\
\midrule\endfirsthead
\toprule
Identifiant & Libellé à l'écran & Type apparent & Valeur observée \\
\midrule\endhead
\multicolumn{4}{@{}l}{\textbf{Onglets}} \\
\texttt{REL-PAT.T01} & Général & onglet & — \\
\texttt{REL-PAT.T02} & Contacts & onglet & — \\
\texttt{REL-PAT.T03} & Quota soins primaires & onglet & — \\
\addlinespace
\multicolumn{4}{@{}l}{\textbf{Données principales}} \\
\texttt{REL-PAT.D01} & N° IPP & texte & — \\
\texttt{REL-PAT.D02} & Nom & texte & — \\
\texttt{REL-PAT.D03} & Nom de famille 1 & texte & — \\
\texttt{REL-PAT.D04} & Nom de famille 2 & texte & — \\
\texttt{REL-PAT.D05} & Sexe & liste & — \\
\texttt{REL-PAT.D06} & D. Nai & date & — \\
\texttt{REL-PAT.D07} & Âge & entier & — \\
\texttt{REL-PAT.D08} & Num & non\ determine & — \\
\texttt{REL-PAT.D09} & Type pièce d'identité & liste & \texttt{C} / \texttt{C.I.N} \\
\texttt{REL-PAT.D10} & N° pièce d'identité & texte & — \\
\texttt{REL-PAT.D11} & E. Civil & liste & \texttt{C} / \texttt{CÉLIBATAIRE} \\
\texttt{REL-PAT.D12} & Type patient & liste & — \\
\texttt{REL-PAT.D13} & Date photo & date & — \\
\texttt{REL-PAT.D14} & Modifié par & texte & — \\
\texttt{REL-PAT.D15} & Créé par & texte & — \\
\texttt{REL-PAT.D16} & Date d'attribution & date & — \\
\addlinespace
\multicolumn{4}{@{}l}{\textbf{Compagnie d'assurance}} \\
\texttt{REL-PAT.A01} & Compagnie d'assur. & liste & \texttt{00116} / \texttt{CNSS PIPC} \\
\texttt{REL-PAT.A02} & Police & texte & — \\
\texttt{REL-PAT.A03} & N° Assu & texte & — \\
\texttt{REL-PAT.A04} & Profession & texte & — \\
\texttt{REL-PAT.A05} & Num. inscription & texte & — \\
\texttt{REL-PAT.A06} & Date inscription & date & — \\
\addlinespace
\multicolumn{4}{@{}l}{\textbf{Domicile}} \\
\texttt{REL-PAT.H01} & Type & liste & — \\
\texttt{REL-PAT.H02} & Adresse & texte & — \\
\texttt{REL-PAT.H03} & Code postal & texte & — \\
\texttt{REL-PAT.H04} & État & liste & — \\
\texttt{REL-PAT.H05} & Ville & liste & — \\
\texttt{REL-PAT.H06} & Quartier & texte & — \\
\texttt{REL-PAT.H07} & Nationalité & liste & \texttt{504} / \texttt{MAROC} \\
\texttt{REL-PAT.H08} & Téléphone 1 & texte & — \\
\texttt{REL-PAT.H09} & Téléphone 2 & texte & — \\
\texttt{REL-PAT.H10} & Téléphone 3 & texte & — \\
\texttt{REL-PAT.H11} & Téléphone 4 & texte & — \\
\texttt{REL-PAT.H12} & Avertissements SMS & case\ a\ cocher & — \\
\texttt{REL-PAT.H13} & E-mail & texte & — \\
\texttt{REL-PAT.H14} & Avertissements e-mail & case\ a\ cocher & — \\
\texttt{REL-PAT.H15} & Environnement & liste & — \\
\addlinespace
\multicolumn{4}{@{}l}{\textbf{Né}} \\
\texttt{REL-PAT.N01} & Nom. Père & texte & — \\
\texttt{REL-PAT.N02} & Nom. Mère & texte & — \\
\texttt{REL-PAT.N03} & Lieu de naissance - État & liste & — \\
\texttt{REL-PAT.N04} & Ville & liste & — \\
\texttt{REL-PAT.N05} & Pays & liste & \texttt{504} / \texttt{MAROC} \\
\texttt{REL-PAT.N06} & Quartier & texte & — \\
\addlinespace
\multicolumn{4}{@{}l}{\textbf{Commentaire}} \\
\texttt{REL-PAT.K01} & Commentaire & zone\ texte & — \\
\addlinespace
\bottomrule
\end{longtable}

\section{Donner rendez-vous}

Quarante-deux éléments, dont sept boutons de barre d'outils hors tableau. Les trente-cinq
restants se répartissent en quatre blocs.

\begin{longtable}{@{}p{2.9cm}p{3.9cm}p{2.8cm}p{4.0cm}@{}}
\toprule
Identifiant & Libellé à l'écran & Type apparent & Valeur observée \\
\midrule\endfirsthead
\toprule
Identifiant & Libellé à l'écran & Type apparent & Valeur observée \\
\midrule\endhead
\multicolumn{4}{@{}l}{\textbf{Identification du patient}} \\
\texttt{REL-RDV.I01} & N° IPP & texte & — \\
\texttt{REL-RDV.I02} & Téléphone & texte & — \\
\texttt{REL-RDV.I03} & Compagnie & liste & — \\
\texttt{REL-RDV.I04} & N° Assu & texte & — \\
\texttt{REL-RDV.I05} & Âge & entier & — \\
\texttt{REL-RDV.I06} & Date naissance & date & — \\
\texttt{REL-RDV.I07} & C.I.N & texte & — \\
\texttt{REL-RDV.I08} & Num & non\ determine & — \\
\texttt{REL-RDV.I09} & Adresse & texte & — \\
\texttt{REL-RDV.I10} & Ville & liste & — \\
\addlinespace
\multicolumn{4}{@{}l}{\textbf{Rendez-vous}} \\
\texttt{REL-RDV.R01} & Agenda & liste & — \\
\texttt{REL-RDV.R02} & Activité & liste & — \\
\texttt{REL-RDV.R03} & Origine & liste & \texttt{AU} / \texttt{AUTRES HÔPITAUX} \\
\texttt{REL-RDV.R04} & Hôpital/C.S. & liste & — \\
\texttt{REL-RDV.R05} & Médecin ext. & texte & — \\
\texttt{REL-RDV.R06} & Service ext. & texte & — \\
\texttt{REL-RDV.R07} & Observations & zone\ texte & — \\
\texttt{REL-RDV.R08} & Date rendez-vous & horodatage & \texttt{07/28/2026 10:57:07 AM} \\
\texttt{REL-RDV.R09} & Rendez-vous supplémentaire & case\ a\ cocher & — \\
\texttt{REL-RDV.R10} & Type d'attention & liste & — \\
\texttt{REL-RDV.R11} & État & liste & \texttt{En instance} \\
\texttt{REL-RDV.R12} & Durée & entier & \texttt{0} \\
\texttt{REL-RDV.R13} & Date réception & horodatage & — \\
\texttt{REL-RDV.R14} & Imprimer données & case\ a\ cocher & — \\
\addlinespace
\multicolumn{4}{@{}l}{\textbf{Contrôle de modifications}} \\
\texttt{REL-RDV.C01} & Créé par & texte & — \\
\texttt{REL-RDV.C02} & Date création & horodatage & \texttt{07/28/2026 10:57:06 AM} \\
\texttt{REL-RDV.C03} & Modifié par & texte & — \\
\texttt{REL-RDV.C04} & Date mod. & horodatage & — \\
\texttt{REL-RDV.C05} & Confirmé par & texte & — \\
\texttt{REL-RDV.C06} & Date conf. & horodatage & — \\
\texttt{REL-RDV.C07} & Annulé par & texte & — \\
\texttt{REL-RDV.C08} & Date annul. & horodatage & — \\
\addlinespace
\multicolumn{4}{@{}l}{\textbf{Liste d'attente des consultations}} \\
\texttt{REL-RDV.L01} & Service & colonne & — \\
\texttt{REL-RDV.L02} & Agenda & colonne & — \\
\texttt{REL-RDV.L03} & Activité & colonne & — \\
\addlinespace
\bottomrule
\end{longtable}




\end{document}
